\worldchapter{Zauber}{appendix-spells}
\label[appendix]{app:spells}

\begin{multicols*}{2}

\section{Astrale Magie}

\subsection{Kristallbarriere}

\spellheader{Beschwörung}{Erde}{1 Kampfhandlung}
An einer Stelle in Reichweite entsteht eine Barriere aus gewachsenen Kristallen, die recht hart und daher schwer zu überwinden ist. Die Kristalle sehen Bergkristallen gleich, die wie normale Kristalle aus dem Boden wachsen. Die Kristallbarriere kann eine maximale Tiefe von \textbf{Stärke} cm und eine Breite von \textbf{Stärke} Metern haben. Es dauert etwa \textbf{Magieniveau}×5 Runden, sich durch die Barriere zu schlagen.

\spellinfobox{5}{1}{5 Schritt}{---}{Wand}{1 Kampfhandlung}{Nein}


\subsection{Licht der Sterne}

\spellheader{Verzauberung}{Licht}{1 Kampfhandlung}
Nach ca. 20 Minuten beginnen die Augen des Zaubernden matt zu werden, und er sieht sowohl am Tage als auch in der Nacht. Helles Licht blendet den Zaubernden nicht. Bei völliger Dunkelheit kann der Zaubernde jedoch nicht sehen. Der Zauber dauert \textbf{Stärke}+\textbf{Magieniveau} Stunden an.

\spellinfobox{5}{1}{0 Schritt}{Stärke+Magieniveau Stunden}{---}{1 Kampfhandlung}{Nein}


\subsection{Weg der Sterne}

\spellheader{Weissagung}{Licht}{1 Kampfhandlung}
Ein heller Stern weist dem Zaubernden den Weg an sein angedachtes Ziel.

\spellinfobox{3}{1}{0 Schritt}{Stärke+Magieniveau Nächte}{---}{1 Kampfhandlung}{Nein}


\subsection{Sternenbotschaft}

\spellheader{Transmutation}{Licht}{1 Kampfhandlung}
Der Astrologe betrachtet den Sternenhimmel und murmelt wiederholend die zu übermittelnde Botschaft.

Nach ca. 35 Minuten beginnen einzelne Sterne der Gestirne heller zu leuchten als andere. Geübte Astrologen und Astralmagiker können aus diesen Konstellationen eine ca. ein \textbf{Stärke des Zaubers} Sätze lange Nachricht herauslesen, egal wo sie auf der Welt gerade stehen.

\spellinfobox{3}{1}{0 Schritt}{Magieniveau Nächte}{---}{1 Kampfhandlung}{Nein}


\subsection{Große Magieanalyse}

\spellheader{Weissagung}{Arkana}{1 Kampfhandlung}
Der Zaubernde ist in der Lage eine Analyse eines gewirkten oder gerade im wirken begriffenen Zaubers durchzuführen. Der Zaubernde erkennt die Schule der Magie, Wesen des Zaubers (ob nun Heilung, Schaden, Art des Elements, Dauer usw.) sowie eine grobe Einschätzung der Stärke des Zaubers.

\spellinfobox{5}{1}{60 Schritt}{---}{Sphäre}{1 Kampfhandlung}{Nein}


\subsection{Stille des Kosmos}

\spellheader{Beschwörung}{Licht}{1 Kampfhandlung}
Der Zaubernde beschwört die Stille des Kosmos. Diese Stille lässt sowohl die magischen Handlungen als auch die Geräusche verstummen. Der Bereich der Stille hat einen Durchmesser von \textbf{Magieniveau} Schritt und bewegt sich mit dem Zaubernden.

Der Mindestwurf für Zauber und Magiekunde ist im Kreis um 4 erhöht, dies gilt auch für den Zaubernden. Alle Geräusche werden von der Stille verschluckt. Der Zaubernde hört für die Dauer des Zaubers nichts.

Der Zauber dauert \textbf{Stärke} Minuten.

\spellinfobox{8}{2}{0 Schritt}{Stärke Minuten}{Kreis}{1 Kampfhandlung}{Nein}


\subsection{Bestrahlen}

\spellheader{Schaden}{Licht}{1 Kampfhandlung}
Der Zaubernde starrt das Ziel an und beschwört das Licht der Sonne. Aus seinen Augen kommen helle Strahlen, die das Ziel mitten im Gesicht treffen.

Die Strahlen blenden das Opfer, es kann kaum noch etwas sehen. Alle Wahrnehmungswürfe und Angriffe haben einen um \textbf{Magieniveau} erhöhten Mindestwurf.

Der Zauber bleibt für \textbf{Stärke} Kampfrunden aktiv.

\spellinfobox{5}{1}{0 Schritt}{Stärke Runden}{Strahl}{1 Kampfhandlung}{Ja}


\section{Chimärologie}

\subsection{Fähigkeit übernehmen}

\spellheader{Kontrolle}{Natur}{1 Kampfhandlung}
Der Zaubernde berührt ein Tier und beginnt damit den Vers wiederholend zu murmeln.

Bei gelungenem Zauber erhält der Begabte für \textbf{Stärke} Stunden die Fähigkeit des Tieres. Ihm wachsen z.B. flugfähige Flügel, oder er erhält die Nase eines Hundes. Die damit verbundene Transformation kann sich über mehrere Minuten hinziehen und mitunter äusserst schmerzhaft sein.

Der Zaubernde erhält auf ein dem Tier entsprechendes Attribut einen Bonus von \textbf{Magieniveau}.

\spellinfobox{8}{2}{0 Schritt}{Stärke Stunden}{---}{1 Kampfhandlung}{Nein}


\subsection{Leviathanschaffung}

\spellheader{Transmutation}{Natur}{1 Kampfhandlung}
Die Eier eines Krokodils werden zuvor in schwarze Tücher gewickelt.

Der Zaubernde umschließt ein oder mehrere Krokodileier mit seinen Händen und spricht auf sie die Verse der Leviathanschaffung. Anschließend werden die Eier in dunkler und warmer Umgebung bis zum Schlüpfen aufbewahrt.

Aus den Eiern schlüpfen in 10-\textbf{Stärke} Monaten \textbf{Magieniveau} kleine Leviathane.

\spellinfobox{12}{4}{0 Schritt}{10-Stärke Monate}{---}{1 Kampfhandlung}{Ja}


\subsection{Unheilige Bindung}

\spellheader{Transmutation}{Natur}{1 Kampfhandlung}
Der Zaubernde presst alle Tiere, die er vereinen will, für die ganze Handlunsgdauer fest zusammen und wirkt den Zauber.

Die Tiere vereinen sich für \textbf{Stärke}×10 Minuten nach der Vorgabe des Zaubernden zu einer Chimäre. Nach Ablauf oder Abbruch des Zaubers trennen sie sich wieder zu ihrer ursprünglichen Form. Die Chimäre ist agressiv und nicht unter der Kontrolle des Zaubernden.

Die Chimäre darf maximal aus \textbf{Magieniveau} Tieren bestehen.

\spellinfobox{12}{3}{0 Schritt}{---}{---}{1 Kampfhandlung}{Nein}


\section{Dämonologie}

\subsection{Schwarzer Ruf}

\spellheader{Beschwörung}{Dämonisch}{Ritual}
Der Zaubernde ruft die Erscheinung eines niederen Dämonen in die Welt. Der Diener erscheint innerhalb der nächsten 3W6 Minuten. Es erfolgt keine Bindung, das Wesen hat seinen eigenen Willen. Für \textbf{Stärke}×10 Minuten bleibt der Dämon auf der Welt.

Der Mindestwurf dieses Zaubers ist 7-\textbf{Magieniveau}. Modifikatoren des Charakters werden danach abgezogen.

\spellinfobox{15}{5}{0 Schritt}{Stärke×10 Minuten}{---}{Ritual}{Nein}


\subsection{Unnatürlicher Wuchs}

\spellheader{Transmutation}{Dämonisch}{1 Kampfhandlung}
Der Zäubernde kauert sich zusammen und schneidet sich mit einem Ritualdolch in das Fleisch, so dass Blut auf den Boden tropft. Er schließt die Augen und stellt sich den Wuchs vor.

Aus dem Körper des Zaubernden wächst die von ihm vorgestellte Form. Die Oberfläche und das Aussehen richten sich dabei nach der Erscheinung des ihm zugetanen Erzdämons oder seiner Diener. Der Zaubernde kann so ein nahezu beliebig geformtes Körperteil schaffen, welches sich nahezu beliebig bewegen lässt. Das Körperteil kann eine Länge von \textbf{Stärke} Metern haben.

\spellinfobox{9}{2}{0 Schritt}{Magieniveau Stunden}{---}{1 Kampfhandlung}{Ja}


\subsection{Tödlicher Stein}

\spellheader{Schaden}{Dämonisch}{1 Kampfhandlung}
Der Zaubernde führt einen Angriff mit einer steinernden Waffe. Gewöhnlich wird hierbei ein ritueller Dolch verwendet.

Bei einem gelungenen Angriff zerfliesst die Waffe im Körper des Opfers zu zwei tentakelartigen Auswüchsen aus flüssigem Stein. Die Waffe hat dabei \textbf{Stärke}+\textbf{Magieniveau} Schadenspotential, es dauert allerdings auch 2 Aktionen, die Waffe wieder herauszuziehen (dabei sind die Auswüchse beim Herausziehen bereits wieder verschwunden).

\spellinfobox{7}{2}{0 Schritt}{---}{---}{1 Kampfhandlung}{Nein}


\subsection{Schlund}

\spellheader{Beschwörung}{Dämonisch}{1 Kampfhandlung}
An einer beliebigen Stelle in Reichweite tut sich ein Schlund von \textbf{Magieniveau} Metern Durchmesser und \textbf{Stärke} Metern Tiefe auf.

\spellinfobox{11}{3}{10 Schritt}{5 Minuten}{Kreis}{1 Kampfhandlung}{Nein}


\subsection{Neues Fleisch}

\spellheader{Heilung}{Dämonisch}{1 Kampfhandlung}
Der Zaubernde berührt die Wunde der Zielperson. Er streicht darüber und spricht den Namen des Gönners.

Die Wunde der Zielperson schließt sich komplett. Jeglicher Schaden der in Verbindung mit der Wunde steht wird getilgt. Die Wunde schließt sich mit neuem Fleisch, und der Verwundete fühlt sich sofort wie neu geboren. Das neu entstehende Fleisch ist allerdings dämonischer Natur. Es ist eine undefinierbare Substanz, welche sich ganz natürlich mit dem menschlichen Fleisch verbindet. Niemand kann sagen wie sich das neue Fleisch in Zukunft verhält, ob es vom Körper angenommen wird, oder ob ganz unerwartete Effekte auftreten.

Der Zauber heilt (\textbf{Stärke}+\textbf{Magieniveau})×2 Wunden.

\spellinfobox{11}{3}{0 Schritt}{---}{---}{1 Kampfhandlung}{Nein}


\subsection{Globulus}

\spellheader{Transmutation}{Dämonisch}{1 Kampfhandlung}
Der Zaubernde schafft ein Versteck, indem er eine Blase in der Dämonenwelt schafft, in die er und \textbf{Stärke}×2 weitere Charaktere translokiert werden. Die Blase entsteht irgendwo in der Dämonenwelt, ist durchsichtig und lässt Geruch, jedoch nichts weiteres durch. Nach Beendigung des Zaubers werden die Personen in der Blase zurück translokiert.

Die Globule besteht für \textbf{Magieniveau}×5 Minuten.

\spellinfobox{7}{2}{0 Schritt}{Magieniveau×5 Minuten}{---}{1 Kampfhandlung}{Ja}


\subsection{Dämonische Vision}

\spellheader{Weissagung}{Dämonisch}{1 Kampfhandlung}
Die Sicht des Zaubernden ändert sich, und er empfindet die Welt mit der Sicht eines Dämons. Er erkennt alles Vorkommen dämonischen Ursprungs klar und leuchtend sogar durch Wände, ist allerdings auch durch die Verzerrtheit des Anblickt in gewisser Weise eingeschränkt. Zudem kann der Zaubernde magische Aktivitäten in seinem Sichfeld genau erkennen.

Der Zauber hält \textbf{Stärke}+\textbf{Magieniveau} Runden an.

\spellinfobox{5}{2}{0 Schritt}{Stärke+Magieniveau Runden}{---}{1 Kampfhandlung}{Nein}


\subsection{Atem des Wisgu}

\spellheader{Schaden}{Dämonisch}{1 Kampfhandlung}
Aus dem Mund des Zaubernden werden übel riechender Schleim, Blut und Schmutz geschleudert. Jeder, der in Kontakt mit den Substanzen kommt, wird von einem unnatürlichen Ekel für eine Zeit von \textbf{Stärke}×2 Runden völlig vereinnahmt und erhält den Zustand Geschockt \textbf{Magieniveau}.

Nach der Wirkungszeit bleibt der Schleim bestehen.

\spellinfobox{5}{1}{5 Schritt}{Stärke×2 Runden}{Kegel}{1 Kampfhandlung}{Nein}


\subsection{Bannkreis}

\spellheader{Widerrufung}{Dämonisch}{1 Kampfhandlung}
Der Dämonologe zeichnet mit dem Dolch oder einem anderen Gegenstand ein Pentagramm in einer Umrandung in den Boden oder in die Decke. Das Pentagramm darf maximal einen Durchmesser von \textbf{Stärke} Schritt haben. Je mächtiger der Dämon ist umso komplexer wird die Zeichnung.

Die äußere Umrandung des Pentagramms kann von einem dämonischen Wesen zwar von außen nach innen, jedoch nicht in die entgegengesetzte Richtung durchbrochen werden. Üblicherweise wird ein Bannkreis um einen Beschwörungskreis gezogen um den Dämon bis zum binden an der Stelle zu halten. Es lassen sich jedoch auch unabhängig von Beschwörungen Orte durch den Bannkreis schützen. Innerhalb des Bannkreises hat der Dämon keinerlei magische oder dämonische Kräfte, wohl aber die Fähigkeiten seiner körperlichen Form.

Die folgenden Paraphernalia haben eine Wirkung auf den Mindestwurf des Bannkreis :

\textit{ \textbf{Geeignete Umgebung, Ruhe}: -1 } \textbf{Bannkreis ist aus Blut}: -2 \textit{ \textbf{5 Kerzen}: -1 } \textbf{Je weiterem Dämonologen}: -1 \textit{ \textbf{Die Sterne stehen richtig}: -10 } \textbf{Tieropfer (je)}: -2 \textit{ \textbf{Menschenopfer (je)}: -5 } \textbf{Tempel in der Nähe}: 20 \textit{ \textbf{Tag}: 5 } \textbf{Priester in der Nähe}: 10 \textit{ \textbf{Weihwasser in der Nähe}: 5  } \textbf{Höherer Dämon}: 10 × \textbf{Erzdämon}: 100

Wird der Zauber reversiert so gilt der Bannkreis auch umgekehrt. Er lässt dann keinen Dämon hinein, wohl aber heraus. So kann ein Dämonologe einen zusätzlichen Kreis um sich selbst ziehen.

\spellinfobox{5}{3}{0 Schritt}{Stärke+Magieniveau Nächte}{Kreis}{1 Kampfhandlung}{Nein}


\subsection{Brut}

\spellheader{Transmutation}{Dämonisch}{1 Kampfhandlung}
Der Dämonologe sticht mit seiner rituellen Waffe in den zuvor von ihm gerufenen Dämon.

Der Dämonologe zerteilt das Wesen in \textbf{Stärke} eigenständige Dämonen. Die Dämonen agieren eigenständig und müssen auch eigenständig gebunden (wenn der ursprüngliche Dämon dies noch nicht war) und gebannt werden. Es können nur niedere Wesen geteilt werden, keine Diener oder gar Erzdämonen. Ein Paktierer ist in der Lage die direkten Diener der Erzdämonen zu teilen.

Ab Magieniveau 4+: Es können auch Dämonendiener geteilt werden.

\spellinfobox{5}{1}{1 Schritt}{---}{---}{1 Kampfhandlung}{Nein}


\subsection{Dämonische Form}

\spellheader{Transmutation}{Dämonisch}{1 Kampfhandlung}
Der Dämonologe kauert sich zusammen, schneidet sich mit der Rituellen Waffe ins Fleisch und tropft das Blut auf das dem Erzdämon zugetane Element (Für Nebel reicht Wasser, Magie bedeutet ein Magisches Artefakt).

Der Dämonologe wandelt sich in die Gestalt eines der Diener seines verbundenen Erzdämons. Er übernimmt dabei teilweise die Fähigkeiten des Dämons, wobei der Körper genauso verwundbar ist wie in seiner menschlichen Form. Zudem erhält der Dämonologe nur die körperlichen Fähigkeiten des Dämons bei seiner normalen menschlichen Größe, und keinerlei magische Fähigkeiten.

Die Verwandlung hält \textbf{Stärke} Minuten an. Er erhält einen Bonus von \textbf{Magieniveau} Punkten auf ein passendes Attribut.

\spellinfobox{5}{2}{0 Schritt}{Stärke Minuten}{---}{1 Kampfhandlung}{Nein}


\subsection{Gespinst}

\spellheader{Beschwörung}{Dämonisch}{1 Kampfhandlung}
Der Zaubernde wirft ein Stück Dämonischer Natur (Artefakt, dämonisches Objekt oder Neues Fleisch) an die Stelle, über der die Sphere entstehen soll. Dann wartet er bis er erhört wird.

Tentakel aus dämonischer Substanz wachsen zu einer gitterförmigen Kugel von max. \textbf{Stärke} Schritt. Die Tentakel haben 500 Wunden und sind somit kaum zu durchtrennen. Das Netz kann sowohl Lebewesen ein- als auch aussperren.

\spellinfobox{5}{1}{0 Schritt}{Magieniveau Nächte}{Sphäre}{1 Kampfhandlung}{Nein}


\subsection{Dämon rufen}

\spellheader{Beschwörung}{Dämonisch}{Ritual}
Die Vorschriften für die Anrufung eines Dämonen sind so vielfältig wie zerstritten. Erwiesen ist, dass das Erbringen diverser Paraphernalia der Invokation zuträglich ist. Auch in der Handlung gibt es einige Besonderheiten, die sich positiv auf das Gelingen auswirken. Im Allgemeinen lässt sich sagen, dass die Invokation als Anrufung in geeigneter Athmosphäre stattfinden sollte und ein direktes Rufen des Dämons seitens des Dämonologen ist. Es ist also auch möglich eine Anrufung ohne jegliche Vorbereitung nur mit dem Vers zu machen.

Die folgenden Paraphernalia haben eine Wirkung auf den Mindestwurf der Invokation :

\textit{ \textbf{Magieniveau}: -\textbf{Magieniveau} } \textbf{Geeignete Umgebung, Ruhe}: -1 \textit{ \textbf{Heptagramm gezeichnet}: -1 } \textbf{Heptagramm ist aus Blut}: -2 \textit{ \textbf{Sigille gezeichnet}: -1 } \textbf{Spieler zeichnet Sigille aus der Hand}: -10 \textit{ \textbf{7 Kerzen}: -1 } \textbf{Je weiterem Dämonologen}: -1 \textit{ \textbf{Die Sterne stehen richtig}: -10 } \textbf{Tieropfer (je)}: -4 \textit{ \textbf{Menschenopfer (je)}: -10 } \textbf{Bannkreis gezogen}: obligatorisch \textit{ \textbf{Tempel in der Nähe}: 20 } \textbf{Tag}: 5 \textit{ \textbf{Priester in der Nähe}: 10 } \textbf{Weihwasser in der Nähe}: 5 \textit{ \textbf{Ein Opfer beginnt zu beten}: 2 (je)  } \textbf{Niederer Dämon wird gerufen}: -1 \textit{ \textbf{Höherer Dämon wird gerufen}: 5 } \textbf{Diener eines Erzdämons wird gerufen}: 30 × \textbf{Erzdämon wird gerufen}: 100

Das Rufen eines dämonischen Wesens schließt das Binden des Dämons nicht mit ein.

\spellinfobox{10}{4}{0 Schritt}{---}{---}{Ritual}{Nein}


\subsection{Dämonisches Wesen binden}

\spellheader{Kontrolle}{Dämonisch}{1 Kampfhandlung}
Der Dämonologe hat dem Wesen, dass er binden möchte, in die Augen (so vorhanden) zu blicken, und muss sich seinem Willen stellen.

Gelingt der Zauber, so erlangt der Dämonologe die Kontrolle über einen Dämon. Ist der Dämon ungebunden, so reicht das reine Ausüben des Zaubers zum Binden. Wurde der Dämon jedoch bereits von einem anderen Dämonologen beherrscht, so ist es erforderlich, zunächst (vor Ausübung des Zaubers) einen magischen Vergleich (Vergleichswurf Zaubern) gegen den beherrschenden Zauberer auszuführen. Misslingt dieser, so bleibt der Dämon weiterhin unter der Herrschaft seines Ursprünglichen Herren. Misslingt nach einem gewonnenen magischen Vergleich der Zauber, so ist der Dämon keinem Herren mehr unterworfen.

Der Mindestwurf des Zaubers ist entsprechend des zu Bindenden Wesens verändert:

\textit{ \textbf{Niederer Dämon}: -2 } \textbf{Höherer Dämon}: 2 \textit{ \textbf{Diener eines Erzdämons}: 20 } \textbf{Erzdämon}: 100

Der Mindestwurf wird um das \textbf{Magieniveau} verringert.

\spellinfobox{7}{1}{0 Schritt}{---}{---}{1 Kampfhandlung}{Nein}


\subsection{Pakt}

\spellheader{Transmutation}{Dämonisch}{Ritual}
Nur wenige, die einen Pakt mit einem Erzdämon eingegangen sind, berichteten über den Hergang des Paktes, doch einige Fakten sind bekannt, der Dämonologe muss über einen der Diener kontakt aufnehmen. Dämonologen haben es dabei leichter, da sie in der Lage sind, diese zu beschwören. Andere Ausrichungen müssen sich an eine der Kultstätten des Dämons begeben.

Hat ein Zaubernder kontakt aufgenommen übernimmt der Dämon die Kontrolle, zumeist wird zu diesem Zwecke vom Dämon ein Portal in seine Globule der Dämonensphäre geöffnet, nur jene, die diese Schwelle übertreten können ohne zu vergehen, haben überhaupt die Aussicht auf einen Pakt.

Was genau in der Dämonensphäre geschieht ist ungewiss, jedoch wird von grausamen Prüfungen berichtet die der Begabte zu überstehn hat, selbst die Stärksten kehren zumeist gebrochen zurück.

Sollte der Pakt erfolgreich geschlossen werden, kehren die Dämonologen als andere Wesen zurück. Zumeist erinnert nur die Erscheinung an den, der in das Portal trat.

In jedem Fall sind die Paktierer nun Untergebene des Dämons, Ungehorsamkeiit wird bestraft, sofort und von innen heraus, dabei spielt es keine Rolle wo sich der Paktierer befindet. Der Pakt verbindet Dämon und Dämonologen über alle Sphären und Unpässlichkeiten hinweg.

Einige besonders mächtige Dämonologen sind zu Beginn noch in der Lage zu wiederstehen und gar ihren eigenen Willen gegen den Dämon durchzusetzen, dennoch sind alle Paktierer früher oder später ihrem Herren untergeben.

Mit dem schliessen eines Paktes wird nicht nur das Leben dem Dämon verschrieben, sondern auch jegliche Existenz nach dem Tod, der Dämonologe stirbt erst dann, wenn der Dämon es zulässt ansonsten wird er lediglich in die Dämonensphäre gezogen um vom Dämon nach belieben wieder entlassen zu werden.

Angeblich solll es Dämonen geben, die Paktierern ihren Pakt entziehen, das wäre die einzige Möglichkeit für einen Paktierer wieder ein halbwegs normales Leben anzutreten, jedoch ziehen es die meisten es vor den Dämonologen zu töten oder zu einem niederen Untergebenen zu machen, sollte er sich wiedersetzen.

\spellinfobox{30}{12}{0 Schritt}{---}{---}{Ritual}{Nein}


\subsection{Mephitische Wolke}

\spellheader{Beschwörung}{Dämonisch}{1 Kampfhandlung}
Beschwört eine Wolke aus giftigen Gasen, die \textbf{Magieniveau} W6 Kampfrunden lang bestehen bleibt. Die Wolke hat einen Durchmesser von \textbf{Stärke} Schritt, und kann bis zu 15 Schritt entfernt vom Zaubernden beschworen werden.

Beendet ein Charakter seine Kampfrunde in der Wolke, erhält er "Vergiftet 2" und 2 Wunden. Durchquert ein Charakter die Wolke, ohne seine Kampfrunde darin zu beenden, erhält er "Vergiftet 1".

\spellinfobox{8}{2}{15 Schritt}{Magieniveau W6 Runden}{Wolke}{1 Kampfhandlung}{Nein}


\subsection{Schattenriss}

\spellheader{Schaden}{Dämonisch}{1 Kampfhandlung}
Der "Schattenriss" ist ein mächtiger und grausamer Zauber, der die dunklen Kräfte der Dämonensphäre beschwört, um die Gelenke eines Gegners mit unvorstellbarer Gewalt auseinanderzureißen. Der Zaubernden beschwört dunkle, tentakelartige Schatten, die sich um die Gliedmaßen des Ziels wickeln und mit einem schaurigen Knacken die Gelenke auseinanderziehen. Im schlimmsten Fall kann dies zum vollständigen Verlust der betroffenen Extremität führen. Sollte das Ziel an den betreffenden Gelenken spezielle Rüstung wie Arm oder Beinschienen tragen so wird die Anzahl der Schutzpunkte von den Erfolgen abgezogen.

Ziel: Ein einzelnes Lebewesen innerhalb der Sichtweite des Zaubernden.

Wirkung: Das Ziel erleidet schweren Schaden an den Gelenken, was zu erheblichen Bewegungseinschränkungen führt. Bei einem besonders mächtigen Wurf kann eine Extremität vollständig abgetrennt werden. Der Schaden beträgt Stärke plus Magieniveau.

Dauer: Sofortiger Effekt, mit anhaltenden Bewegungseinschränkungen, bis das Ziel geheilt wird.

Nebenwirkungen: Der Einsatz dieses Zaubers kann die Aufmerksamkeit dunkler Mächte erregen, die den Zaubernden in Zukunft heimsuchen könnten.

\begin{pquote}{Totschläger-Arne, Räuber}
"Als wir in den Wald stürmten stand da nur diese junge Frau. Unser Hauptmann lachte. Das sollte alles sein, was uns diese Menschen hier entgegen werfen? Eine Hexe? Unser Hauptmann stürmte ihr entgegen, doch die Hexe erhob nur ihre Hände. Im gleichen Moment wurde der erhobene Schwertarm unseres Hauptmanns von dunklen nebligen Tentakeln zurück gerissen mit einem gräßlichen Geräusch. Er schrie. Als die Hexe erneut ihre dämonische Geste vollführte, kamen weitere Tentakel und rissen ihm den ganzen Arm am Gelenk aus! Wir konnten grade noch fliehen!"
\end{pquote}

\spellinfobox{5}{2}{10 Schritt}{---}{Strahl}{1 Kampfhandlung}{Nein}


\section{Echsenmagie}

\subsection{Blutopfer-Ritual}

\spellheader{Heilung}{Blut}{Ritual}
In einem grausamen Ritual opfert der Zaubernde ein kleines Lebewesen, um vorübergehend Boost auf ein beliebiges Attribut in Höhe der \textbf{Zauberstärke} zu erhalten. Der Boost hält für \textbf{Magieniveau} Stunden an.

\spellinfobox{10}{3}{0 Schritt}{Magieniveau Stunden}{---}{Ritual}{Nein}


\subsection{Geistesbann}

\spellheader{Kontrolle}{Geist}{1 Kampfhandlung}
Das Ziel würfelt mit seinem Logikwert, der Wurf ist um die \textbf{Zauberstärke} erschwert.

Misslingt der Wurf, ist das Ziel gezwungen, den Befehlen des Zaubernden zu gehorchen, bis der Zauber abläuft oder aufgehoben wird.

\spellinfobox{12}{3}{10 Schritt}{Magieniveau×5 Stunden}{---}{1 Kampfhandlung}{Nein}


\subsection{Säurehauch}

\spellheader{Schaden}{Dämonisch}{1 Kampfhandlung}
Der Zaubernde spuckt eine Säurewolke, die allen Zielen im Wirkungsbereich Schaden in Höhe der \textbf{Zauberstärke}+\textbf{Magieniveau} zufügt.

\spellinfobox{8}{2}{5 Schritt}{---}{Wolke}{1 Kampfhandlung}{Nein}


\subsection{Schuppenrüstung}

\spellheader{Transmutation}{Blut}{1 Kampfhandlung}
Der Körper des Zaubernden wird mit zusätzlichen Schuppen bedeckt. Er erhält \textbf{Zauberstärke} Normalen Schutz und \textbf{Magieniveau} Blutungsschutz.

\spellinfobox{6}{3}{0 Schritt}{---}{---}{1 Kampfhandlung}{Nein}


\subsection{Giftiger Biss}

\spellheader{Schaden}{Blut}{1 Kampfhandlung}
Der Zaubernde erhält für \textbf{Magieniveau} Kampfrunden einen Giftbiss, der bei Treffern zusätzlichen Giftschaden in Höhe der \textbf{Zauberstärke} verursacht.

\spellinfobox{7}{3}{0 Schritt}{3 Runden}{---}{1 Kampfhandlung}{Nein}


\subsection{Dunkelsicht}

\spellheader{Transmutation}{Blut}{1 Kampfhandlung}
Der Zaubernde kann in völliger Dunkelheit bis zur \textbf{Zauberstärke} × \textbf{Magieniveau} Meter sehen.

\spellinfobox{5}{1}{10 Schritt}{8 Stunden}{---}{1 Kampfhandlung}{Nein}


\subsection{Schattenschritt}

\spellheader{Illusion}{Licht}{1 Kampfhandlung}
Der Zaubernde verschmilzt mit den Schatten und erhöht seine Heimlichkeit um die \textbf{Zauberstärke}.

\spellinfobox{6}{2}{0 Schritt}{Magieniveau Minuten}{---}{1 Kampfhandlung}{Nein}


\subsection{Sumpffieber}

\spellheader{Verzauberung}{Natur}{Ritual}
Das Ziel erleidet an jedem Tag, an dem der Fluch aktiv ist, Schaden in Höhe der \textbf{Zauberstärke}. Zusätzlich werden alle körperlichen Attribute um die \textbf{Magieniveau} verringert.

\spellinfobox{12}{4}{0 Schritt}{1 Wochen}{---}{Ritual}{Nein}


\subsection{Echsenruf}

\spellheader{Beschwörung}{Blut}{1 Kampfhandlung}
Ruft eine Anzahl kleiner Echsen herbei, die einfache Befehle ausführen. Die Anzahl entspricht der \textbf{Stärke} des Zaubers.

\spellinfobox{8}{3}{0 Schritt}{Magieniveau Stunden}{---}{1 Kampfhandlung}{Nein}


\subsection{Verwesung}

\spellheader{Schaden}{Blut}{1 Kampfhandlung}
Lässt das Fleisch des Ziels verrotten, was direkten Schaden und Attributsverlust entsprechend der \textbf{Zauberstärke}+\textbf{Magieniveau} verursacht. Das Ziel des Zaubers wählt die Attribute und verteilt die Minuspunkte darauf.

Die Mali auf den Attributen werden erst entfernt, wenn der Schaden vollständig geheilt ist.

\spellinfobox{10}{4}{15 Schritt}{---}{---}{1 Kampfhandlung}{Nein}


\subsection{Schlangenblick}

\spellheader{Kontrolle}{Licht}{1 Kampfhandlung}
Der Zaubernde fixiert das Ziel mit einem hypnotischen Blick und lähmt es für Runden entsprechend der \textbf{Zauberstärke}+\textbf{Magieniveau}.

\spellinfobox{9}{3}{10 Schritt}{Zauberstärke Runden}{---}{1 Kampfhandlung}{Ja}


\subsection{Säuredorn}

\spellheader{Schaden}{Blut}{1 Kampfhandlung}
Ein scharfer Säurendorn wird auf das Ziel geschleudert und verursacht Giftschaden entsprechend der \textbf{Zauberstärke}. Der Dorn hat einen Durchschlag von \textbf{Magieniveau}.

\spellinfobox{6}{2}{15 Schritt}{---}{---}{1 Kampfhandlung}{Nein}


\subsection{Kriechende Furcht}

\spellheader{Kontrolle}{Geist}{1 Kampfhandlung}
Alle Lebewesen in Reichweite des Zaubers würfeln auf ihre Willenskraft, der Wurf ist um die \textbf{Zauberstärke} erschwert.

Wer den Wurf nicht schafft, flieht panisch vor der Echse.

\spellinfobox{10}{4}{15 Schritt}{Magieniveau Minuten}{---}{1 Kampfhandlung}{Ja}


\section{Elementarmagie}

\subsection{Energieblitz}

\spellheader{Schaden}{Energie}{1 Kampfhandlung}
Der Zaubernde beschwört einen Blitz aus reiner Energie, welcher auf ein Ziel zu fliegt und beim Einschlag Treffer in der Höhe der \textbf{Stärke des Zaubers}+\textbf{Magieniveau} verursacht.

\spellinfobox{5}{1}{5 Schritt}{---}{---}{1 Kampfhandlung}{Nein}


\subsection{Flammender Tod}

\spellheader{Schaden}{Feuer}{1 Kampfhandlung}
In einer Entfernung von maximal 10 Metern entsteht ein loderndes Feuer, dass auf einer Fläche von \textbf{Stärke} Metern zum Quadrat  \textbf{Magieniveau} Wunden pro Kampfrunde verursacht. Das Feuer ist nicht magisch und brennt, bis es seine Nahrung verzehrt hat, ohne brennbares Material \textbf{Stärke} Kampfrunden.

\spellinfobox{7}{2}{10 Schritt}{Stärke Runden}{Kreis}{1 Kampfhandlung}{Nein}


\subsection{Eissplitter}

\spellheader{Transmutation}{Wasser}{1 Kampfhandlung}
Der Zaubernde erschafft einen kleinen Eissplitter in seiner Hand, der mit schneller Geschwindigkeit auf das Opfer zufliegt, und es an einer Stelle freier Haut trifft. Der Splitter dringt tief in die Haut ein, schmilzt dann jedoch sofort, so dass nichtmal eine Wunde zurückbleibt

Nach  einer Minute kühlt das getroffene Körperteil soweit ab, dass das Opfer es kaum zu benutzen vermag. Es ist auch kein Gefühl mehr in dem entsprechenden Körperteil vorhanden. Schaden nimmt das Opfer nicht, alle Handlungen mit dem Körperteil verringern die Fertigkeits-/Attributwerte um \textbf{Magieniveau}×2.

Der Effekt dauert \textbf{Stärke}×2 Minuten an.

\spellinfobox{5}{2}{0 Schritt}{Stärke×2 Minuten}{---}{1 Kampfhandlung}{Nein}


\subsection{Wildwuchs}

\spellheader{Beschwörung}{Natur}{1 Kampfhandlung}
Der Zaubernde erzeugt ein unnatürlich schnelles Wachstum von Pflanzen. Innerhalb eines Radius von max. \textbf{Stärke} Meter entstehen natürliche und unnatürliche Pflanzen, die sich durch den Boden graben, durch Mauern dringen und Stahl zum Bersten bringen können. Der Pflanzenwuchs bleibt \textbf{Magieniveau} Tage lang bestehen, danach verrotten die Pflanzen zu einem stinkenden Etwas.

Der Zaubernde kann das Wachstum der Pflanzen nicht kontrollieren.

\spellinfobox{3}{2}{0 Schritt}{Magieniveau Tage}{---}{1 Kampfhandlung}{Nein}


\subsection{Splittersturm}

\spellheader{Schaden}{Erde}{1 Kampfhandlung}
In der Handfläche des Zaubernden bilden sich Splitter, die mit hoher Geschwindigkeit auf das Ziel zufliegen. In einem Winkel von \textbf{Magieniveau}×15° verursachen die Splitter insgesamt \textbf{Stärke}×3 Treffer, und verursachen große Strukturschäden an festen Objekten.

Werden Lebewesen getroffen, verteilt der Spielleiter den Schaden auf die Opfer.

\spellinfobox{9}{3}{10 Schritt}{---}{Kegel}{1 Kampfhandlung}{Nein}


\subsection{Sichtschirm}

\spellheader{Beschwörung}{Luft}{1 Kampfhandlung}
Der Zaubernde verschwimmt vor dem Hintergrund. Einzig ein leichtes Flackern in der Luft verrät das Vorhandensein eines Objektes an der Stelle. Bewegt sich der Zaubernde, so bewegt sich auch das Sichtschild mit. Für das Entdecken des ungewöhnlichen Schimmerns ist ein Wahrnehmungswurf notwendig, der Erfolge gemäß der \textbf{Stärke} des Zaubers erreicht.

\spellinfobox{5}{2}{0 Schritt}{Magieniveau Runden}{---}{1 Kampfhandlung}{Nein}


\subsection{Schmelze}

\spellheader{Transmutation}{Natur}{1 Kampfhandlung}
Auf einer Fläche von \textbf{Stärke}W6 Quadratmetern in 2 Schritt Entfernung verflüssigt sich der Untergrund. Nach \textbf{Magieniveau} Minuten erstarrt der Boden  innerhalb von drei Sekunden wieder.

\spellinfobox{7}{3}{2 Schritt}{Magieniveau Minuten}{---}{1 Kampfhandlung}{Nein}


\subsection{Odem}

\spellheader{Beschwörung}{Dämonisch}{1 Kampfhandlung}
Für \textbf{Stärke} Kampfrunden ist der Atem des Zaubernden ein fauliger, dämonischer Schwall, welcher eine Reichweite von 2 Metern hat und jedem innerhalb der Wolke pro Kampfrunde \textbf{Magieniveau} Wunden zufügt.

\spellinfobox{5}{1}{2 Schritt}{Stärke Runden}{---}{1 Kampfhandlung}{Nein}


\subsection{Nebelschleier}

\spellheader{Beschwörung}{Wasser}{1 Kampfhandlung}
Aus dem Boden um den Zaubernden steigt, während er die Arme hebt, ein dichter Nebel auf, der die Sicht hemmt und Geräusche dämpft. Die Nebelwolke hat einen Durchmesser von \textbf{Stärke}×3 Metern und eine Höhe von etwa 3 Metern. Je höher die \textbf{Stärke} des Zaubers ist, desto dichter ist auch die Nebelwolke. Die Wolke ist stationär und bildet sich um den Zaubernden als Zentrum.

Der Nebel wirkt durch die magische Verbindung zum Zaubernden wie eine Erweiterung seiner Sinne. Solange er sich selbst im Nebel aufhält, kann er alle Bewegungen innerhalb der Wolke instinktiv wahrnehmen und alle Geräusche in ihr besser hören (Wahrnehmung + \textbf{Magieniveau}).

Es hält sich das Gerücht, daß lautes, manisches Gelächter bei der Beschwörung des Nebels dessen spätere, bedrohliche Wirkung erhöht. Dies kann jedoch mit ziemlicher Sicherheit in das Reich der Mythen und Legenden abgetan werden.

Die Nebelwolke bleibt für \textbf{Stärke} Minuten bestehen.

\spellinfobox{5}{1}{0 Schritt}{Stärke Minuten}{---}{1 Kampfhandlung}{Nein}


\subsection{Manasturm}

\spellheader{Beschwörung}{Energie}{1 Kampfhandlung}
Kurze Zeit nach der Anrufung manifestiert sich direkt über dem Zaubernden ein magischer Nebel, der Blitze absondert und eine magische Spannung erzeugt. Der Sturm wächst mit einer Geschwindigkeit von einem Meter pro Kampfrunde bis zu einer Größe von \textbf{Stärke}×5 Metern an, und lässt sich durch den Zaubernden steuern.

Im magischen Sturm erleidet jeder Magiekundige eine Wunde pro Kampfrunde. Zudem ist kein Begabter im Sturm in der Lage, eine magische Handlung zu vollziehen. Nicht Magiekundige erleiden keine Einschränkungen.

Der Sturm lässt sich mit einer Geschwindigkeit von zwei Metern pro Aktion steuern. Dieses Steuern erfordert, dass sich der Zaubernde weiter auf den Zauber konzentriert, was ansonsten nicht von Nöten ist.

\spellinfobox{7}{2}{0 Schritt}{Stärke Runden}{---}{1 Kampfhandlung}{Nein}


\subsection{Lebensstrom}

\spellheader{Heilung}{Wasser}{Ritual}
Der Zaubernde entkleidet sich vollkommen und legt sich mit dem Gesicht nach unten in fliessendes Gewässer, welches so gross ist, dass der Zaubernde ganz darin unterkommt. Dort lässt er sich treiben.

Während der ganzen Zeit des Treibens zieht der Zaubernde Lebensenergie dem Fluss, und heilt \textbf{Stärke}+\textbf{Magieniveau} Wunden in einer Stunde. Während dieser Zeit kann und muss er nicht Atmen, und nicht seine Umwelt wahrnehmen.

\spellinfobox{5}{2}{0 Schritt}{---}{---}{Ritual}{Nein}


\subsection{Kohlestein}

\spellheader{Transmutation}{Energie}{1 Kampfhandlung}
Der Zaubernde kann einen Diamanten oder Edelstein von jeder beliebigen Größe in ein glühendes Stück Kohle verwandeln. Das Stück entspricht der Größe des Diamanten und bleibt \textbf{Magieniveau} Stunden am Glühen. Dabei ist es so heiss, dass es brennbare Stoffe entzündet. Je reiner der Edelstein ist, desto heisser glüht das Kohlestück. Ein Bernstein reicht nicht aus, um etwas mit der Kohle zu entzünden.

\spellinfobox{3}{1}{0 Schritt}{Magieniveau Stunden}{---}{1 Kampfhandlung}{Nein}


\subsection{Frischer Wind}

\spellheader{Beschwörung}{Luft}{1 Kampfhandlung}
Frischer Wind lässt einen ermutigenden und frischen Luftzug erscheinen. In Momenten der Hoffnungslosigkeit ist der Zauber das Richtige, um die Unternehmungsfreude der Gefährten zu steigern.

Der Wind streicht etwa eine Minute lang in einem Umkreis von 100 Metern über das Land, und erfüllt alle, die er berührt mit neuem Mut und neuer Frische. Jeder, der sich in dem Radius befindet, heilt \textbf{Stärke} Wunden und erhält für die nächsten zwei Stunden einen Bonus von \textbf{Magieniveau} Punkten auf seine Tapferkeit.

\spellinfobox{3}{1}{0 Schritt}{2 Stunden}{---}{1 Kampfhandlung}{Nein}


\subsection{Fäule}

\spellheader{Schaden}{Dämonisch}{1 Kampfhandlung}
Der Begabte hat auf eine Pflanze oder ein Lebewesen zu spucken.

Solange das Opfer vom Speichel benetzt ist, fault das Fleisch oder die Pflanze immer weiter. Solange das Fleisch fault, verursacht der Zauber kumulativ eine Wunde in jeder 3. Kampfrunde. Endet der Zauber, so fault das Opfer nicht mehr weiter, verdorbenes Fleisch ist jedoch für immer verloren.

Der Zauber endet wenn der Speichel abgewaschen wurde oder nach \textbf{Stärke}+\textbf{Magieniveau} Wunden.

\spellinfobox{9}{3}{3 Schritt}{---}{---}{1 Kampfhandlung}{Nein}


\subsection{Egelranken}

\spellheader{Beschwörung}{Dämonisch}{1 Kampfhandlung}
Der Zaubernde kniet sich auf den Boden fixiert das Ziel mit seinem Blick und schlägt die Finger einer Hand in den Erdboden.

Aus den Fingern des Zaubernden wuchern mit atemberaubender Geschwindigkeit Ranken, die unter der Erdoberfläche in Richtung des Gegners wachsen. Unter der Erde bewegen sich die Ranken mit einer Geschwindigkeit von \textbf{Magieniveau} Metern pro Kampfrunde. Sobald sich die Ranken unter dem Gegner befinden, wachsen sie an die Oberfläche und umschlingen seine Beine, so das dieser weder laufen noch uneingeschränkt Kämpfen kann. Seine Kampffertigkeiten (Nahkampf, Schießen, Werfen) sind um \textbf{Stärke} verringert.

\spellinfobox{5}{1}{15 Schritt}{Stärke Runden}{---}{1 Kampfhandlung}{Nein}


\subsection{Efeumantel}

\spellheader{Beschwörung}{Natur}{1 Kampfhandlung}
Der Zaubernde lässt Efeuranken aus dem Boden wachsen, welche sich um seinen ganzen Körper ranken, und nach dem Wachstum vom Boden lösen, so dass der Zaubernde sich frei bewegen kann.

Die Ranken bieten dem Zaubernden im Kampf einen Schutz von \textbf{Magieniveau} für \textbf{Stärke} Kampfrunden.

Sie verfallen nach einer Stunde zu einem welken Haufen. Bis dahin geben die Ranken auch außerhalb des Kampfes einen Bonus von \textbf{Stärke} auf Heimlichkeit abhängig von der Umgebung.

\spellinfobox{5}{1}{0 Schritt}{1 Stunden}{---}{1 Kampfhandlung}{Nein}


\subsection{Bernsteinpfad}

\spellheader{Transmutation}{Natur}{1 Kampfhandlung}
Der Zaubernde legt an dem Ort, an dem das Bernsteinportal entstehen soll, fünf Bernsteinkristalle in Form eines Pentagramms und konzentiert sich auf den Zielort. Dann stellt er sich eine Minute lang einen Tunnel vor, der ihn zu diesem Ort bringt.

Das Portal bleibt \textbf{Magieniveau} Stunden bestehen und kann \textbf{Stärke}×2 Menschen oder Tiere transportieren.

Am Ort des Bernsteinpentagramms beginnt eine Efeupflanze kreisförmig zu wachsen und einen Vortex zu bilden. Der Vortex bildet in der Mitte einen Schlund welcher durch ein Portal zu einem dem Zaubernden bekannten Ort in einer Entfernung von 20 Meilen führt. Betritt der Zaubernde das Portal, so taucht er sofort auf der anderen Seite am Zielort auf.

\spellinfobox{9}{2}{0 Schritt}{Magieniveau Stunden}{---}{1 Kampfhandlung}{Nein}


\subsection{Auge der Seth'Nra}

\spellheader{Transmutation}{Dämonisch}{1 Kampfhandlung}
Der Zaubernde nimmt mit zwei Fingern eines seiner Augen aus der Augenhöhle und hält es auf der flachen Hand.

Das Auge bekommt schwarze Auswüchse in Form von Tentakeln und Flügeln und beginnt von selbst zu fliegen. Es bewegt sich mit der Geschwindigkeit von 10 Metern pro Sekunde und hat \textbf{Stärke} maximale Wunden. Falls das Auge bei Ablauf des Zaubers nicht zurück ist, fällt es zu Boden. In diesem Fall kann der Zaubernde es innerhalb einer halben Stunde wiederfinden und einsetzen.

Wird das Auge zerstört oder geht es verloren, so verdorrt es und wächst erst nach 2W6 Tagen nach. In diesem Fall verursacht es beim Zaubernden einmalig 2 Wunden.

Der Zauber wirkt für \textbf{Magieniveau}×5 Minuten.

\spellinfobox{7}{2}{0 Schritt}{Magieniveau×5 Minuten}{---}{1 Kampfhandlung}{Nein}


\subsection{Elementare Waffe}

\spellheader{Transmutation}{Energie}{1 Kampfhandlung}
Der Charakter kanalisiert die Magie seines bevorzugten Elements in seine Nahkampfwaffe, um sie zu verstärken.

Für die nächsten \textbf{Magieniveau} Kampfrunden ignoriert die Waffe den Schutz des Gegners. Außerdem haben Angriffe mit der Waffe für die nächsten zwei Kampfrunden ein um \textbf{Stärke} erhöhtes Schadenspotenzial.

\spellinfobox{5}{2}{0 Schritt}{Magieniveau Runden}{---}{1 Kampfhandlung}{Nein}


\subsection{Windpfeil}

\spellheader{Transmutation}{Luft}{1 Kampfhandlung}
Der Zaubernde verzaubert bis zu \textbf{Magieniveau} Pfeile mit dem Element des Winds. Bei einer erfolgreichen Zauberwirkung ignoriert die Waffe Rüstungsschutz. Der Pfeil trifft das Opfer mit einer solchen Kraft das es auf seine Resistenz werfen muss. Erreicht es weniger Erfolge als die \textbf{Stärke} des Zaubers, so wird es zu Boden geworfen.

\spellinfobox{5}{1}{15 Schritt}{---}{---}{1 Kampfhandlung}{Nein}


\subsection{Avatar des Sturms}

\spellheader{Beschwörung}{Luft}{1 Kampfhandlung}
Der Zaubernde verbraucht all sein Arkana und wird zu dem Fokuspunkt eines tosenden Sturms. Für \textbf{Magieniveau} W6 + \textbf{verbrauchtes Arkana} Runden erleiden alle Kreaturen elektrischen Schaden in Höher der \textbf{Stärke des Zaubers}.

Alle Kreaturen im Einflussbereich müssen pro Runde eine Athletik Probe mit 2 Erfolgen bestehen oder werden vom tosenden Wind zu Boden geworfen.

\spellinfobox{15}{3}{25 Schritt}{Magieniveau W6 + Verbrauchtes Arkana Runden}{Wolke}{1 Kampfhandlung}{Nein}


\subsection{Kettenblitz}

\spellheader{Schaden}{Luft}{1 Kampfhandlung}
Der Zaubernde beschwört die Essenz eines Gewitters in seiner Hand und schleudert sie als Blitz auf eine Kreatur. Die getroffene Kreatur erleidet \textbf{Stärke} Treffer.

Der Blitz prallt von ihr ab und bewegt sich zu der Kreatur, die der getroffenen am nächsten ist. Diese erhält \textbf{Stärke}-1 Treffer.

Dieser Vorgang wiederholt sich \textbf{Magieniveau}+2 mal.

Der Zauber unterscheidet nicht zwischen Freund und Feind, und trifft jedes Opfer nur einmal.

\spellinfobox{5}{3}{15 Schritt}{---}{---}{1 Kampfhandlung}{Nein}


\subsection{Flügel des Windes}

\spellheader{Transmutation}{Luft}{1 Kampfhandlung}
Dem Verzauberten wachsen Windflügel, die ihn oder eine andere Kreatur durch die Lüfte tragen.

Für \textbf{Zauberstärke}+\textbf{Magieniveau} Minuten kann die verzauberte Kreatur fliegen.

\spellinfobox{5}{2}{1 Schritt}{Stärke+Magieniveau Minuten}{---}{1 Kampfhandlung}{Nein}


\subsection{Kleiner Sandsturm}

\spellheader{Transmutation}{Luft}{1 Kampfhandlung}
Der Zaubernde konzentriert sich und wirbelt Staub, Erde oder Sand in seiner Sichtweite auf. Es bildet sich ein kleiner Sandsturm von \textbf{Magieniveau} Schritt Höhe und \textbf{Magieniveau} Schritt Breite für \textbf{Stärke} Runden.

Wesenheiten im Zentrum des Sandsturms müssen eine Wahrnehmungsprobe ablegen und brauchen mindestens soviele Erfolge wie die Stärke des Zaubers. Ansonsten sind sie für die Zeit des Zaubers blind mit allen Mali auf Blindheit.

\spellinfobox{5}{1}{10 Schritt}{Stärke Runden}{Wolke}{1 Kampfhandlung}{Nein}


\subsection{Elementar}

\spellheader{Beschwörung}{Luft}{1 Kampfhandlung}
Der Zauberne beschwört für \textbf{Stärke} Runden eine Kreatur aus purer elementarer Energie, welche ihn im Kampf unterstützen kann. Das Elementar hat \textbf{Magieniveau} Wunden und kann mit \textbf{Magieniveau} Schadenspotential angreifen.

\spellinfobox{10}{2}{10 Schritt}{Stärke Runden}{---}{1 Kampfhandlung}{Nein}


\subsection{Wasseratem}

\spellheader{Transmutation}{Wasser}{1 Kampfhandlung}
Der Zaubernde spricht die Worte Þat mælti mín móðir, at mér skyldi kaupa fley ok fagrar árar.

Der Zaubernde verzaubert sich oder eine Persion, die er berühren kann.

Der Verzauberte ist für (\textbf{Zauberstärke}+\textbf{Magieniveau})×2 Minuten in der Lage unter Wasser zu atmen. Nachdem der Zauber endet muss der Verzauberte etwas in seiner natürlichen Sprache sagen, damit er wieder Luft atmen kann.

\spellinfobox{5}{2}{0 Schritt}{(Stärke+Magieniveau)×2 Minuten}{---}{1 Kampfhandlung}{Nein}


\subsection{Elementargeist rufen}

\spellheader{Beschwörung}{Geist}{Ritual}
Der Zaubernde legt ein wenig des Elements, dessen Elementargeist gerufen werden soll in eine Schale oder auf einen Untergrund aus dem gegensätzlichen Element. Dann kniet er vor der Schale nieder und ruft die Kräfte des Elementes an.

Nach etwa 30 Minuten erscheint in dem benutzten Element das Antlitz eines Elementargeistes. Dieser kann je nach Element sehr variieren, so dass von einem Gesicht bis hin zu einer nicht materiellen Erscheinung alles möglich ist. Der Elementargeist hat keine Gefühle und keine Ausrichtung, er lässt sich vom Zaubernden jedoch nur kontrollieren, wenn dieser ihn bindet (Elementargeist binden).

Der Elementargeist bleibt für (\textbf{Stärke}+\textbf{Magieniveau})×2 Minuten.

\spellinfobox{8}{4}{0 Schritt}{(Stärke+Magieniveau)×2 Minuten}{---}{Ritual}{Nein}


\subsection{Elementargeist binden}

\spellheader{Kontrolle}{Geist}{1 Kampfhandlung}
Der Zaubernde deutet auf den Elementargeist.

Der Zaubernde kann einen Elementargeist (welcher sich schon materialisiert haben muss) an sich binden. Ist dies geschehen, folgt der Elementargeist jedem Befehl des Zaubernden. Ist der Elementargeist bereits gebunden, so kann der Zaubernde den Elementargeist lediglich mit einem Magischen Duell übernehmen, indem er dem fremden Zauber übernimmt. Der Zaubernde bindet den Elementargeist für (\textbf{Stärke}+\textbf{Magieniveau})×2 Minuten an sich.

\spellinfobox{5}{2}{0 Schritt}{(Stärke+Magieniveau)×2 Minuten}{---}{1 Kampfhandlung}{Nein}


\subsection{Elementargestalt}

\spellheader{Transmutation}{Natur}{1 Kampfhandlung}
Der Zaubernde berührt das entsprechende Element, konzentriert sich und murmelt "(Element) werde mein Körper".

Der Körper des Zaubernden verwandelt sich in das entsprechende Element, mit allen seinen Vor- und Nachteilen. Es sollte beachtet werden, das Kleidung und Ausrüstung nicht mit verwandelt wird, und gegebenenfalls Schaden nehmen kann. DerZauber kann jederzeit fallengelassen werden. Während der Aufrechterhaltung des Zaubers kann der Zaubernde keinen Elementarzauber wirken der ein anderes Element als das seines Körpers zur Grundlage hat.

Der Zauber hält (\textbf{Stärke}+\textbf{Magieniveau})×2 Minuten an.

\spellinfobox{5}{1}{0 Schritt}{(Stärke+Magieniveau)×2 Minuten}{---}{1 Kampfhandlung}{Nein}


\subsection{Element beschwören}

\spellheader{Beschwörung}{Natur}{1 Kampfhandlung}
Der Zaubernde schliesst die linke Hand zur Faust.

In der Faust des Zaubernden entsteht ein kleines Vorkommen des Elementes, das er herbeirufen möchte. Vorwiegend wird dieser Zauber in Kombination mit dem Rufen eines Elementargeistes genutzt.

\spellinfobox{3}{1}{0 Schritt}{---}{---}{1 Kampfhandlung}{Nein}


\subsection{Elementargegenstand}

\spellheader{Transmutation}{Natur}{1 Kampfhandlung}
Der Zaubernde konzentriet sich auf den gewünschten Gegenstand vor seinem geistigen Auge, greift in ein Objekt aus einem bestimmten Element (es ist ihm bei dem Zauber möglich, seine Hand einfach in das Objekt gleiten zu lassen) und zieht den gewünschten Gegenstand heraus.

Der Zaubernde zieht einen Gegenstand seiner Wahl aus einem anderen Objekt. Denkbar sind Schwerter aus Eis, Trinkgefäße aus Holz, Schilde aus Fels oder Wasser, usw. Dem Zaubernden sind seiner Kreativität keine Grenzen gesetzt, nur kann er immer nur einen Gegenstand herausziehen, was die Herstellung von z.B. Ketten sehr aufwendigt macht.

Der Gegenstand besitzt die für ihn typischen Eigenschaften, plus der elementaren Komponente, so fügt ein Schwert aus Feuer dem Opfer Brandschaden zu, ein Schwert aus Holz erhält keinen Schadensbonus. Nur der Zaubernde erhält keinen diesen möglichen elementaren Schaden. Ein Trinkpokal aus Eis ist sehr angenehm, wenn man gerne sehr stark gekühlte Getränke geniesst, allerdings könnte er allen anderen ausser dem Erschaffer an den Finger und oder Lippen festfrieren.

Der Zauber kann jederzeit vom Zaubernden fallengelassen werden. Bringt der Zaubernde den Gegenstand beim Fallenlassen des Zaubers nicht wieder an seinen Ursprungsort zurück (z.B. steckt er den hölzerner Schild nicht wieder in den Baum zurück) erhält er eine Wunde an der Hand, gegebenenfalls ein paar Brand- oder Frostblasen und der Gegenstand löst sich auf.

Der Gegenstand bleibt (\textbf{Stärke}+\textbf{Magieniveau})×10 Minuten bestehen.

\spellinfobox{5}{2}{0 Schritt}{(Stärke+Magieniveau)×10 Minuten}{---}{1 Kampfhandlung}{Nein}


\subsection{Ginaes Ruf}

\spellheader{Heilung}{Wasser}{1 Kampfhandlung}
Der Zaubernde weist den zu Behandelnden an, sich seiner Kleidung zu entledigen und sich in ein fliessendes Gewässer zu legen. Dort stellt sich der Zaubernde neben den zu Behandelnden und legt seine Hand auf dessen Kopf.

Während der ganzen Zeit des Treibens zieht der zu Behandelnde Lebensenergie entsprechend \textbf{Stärke}+\textbf{Magieniveau} Wunden pro Minute aus dem Fluss. Während dieser Zeit kann und muss er nicht atmen. Er nimmt seine Umwelt nicht wahr. Der Zaubernde muss während der gesamten Zeit neben ihm stehen.

\spellinfobox{2}{1}{0 Schritt}{---}{---}{1 Kampfhandlung}{Nein}


\subsection{Kälte}

\spellheader{Transmutation}{Eis}{1 Kampfhandlung}
Der Zaubernde richtet seine Hand auf das Opfer des Zaubers. Dann spreizt er die Finger und ruft Jogran an.

Der Zaubernde erschafft einen kleinen Eissplitter in seiner Hand, der mit schneller Geschwindigkeit auf das Opfer zufliegt, und es an einer Stelle freier Haut trifft. Der Splitter dringt tief in die Haut ein, schmilzt dann jedoch sofort, so dass nichtmal eine Wunde zurückbleibt

In den nächsten 5 Minuten kühlt das getroffene Körperteil soweit ab, dass das Opfer es kaum zu benutzen vermag. Es ist auch kein Gefühl mehr in dem entsprechenden Körperteil vorhanden. Schaden nimmt das Opfer nicht, alle Handlungen mit dem Körperteil verringern die Fertigkeits-/Attributwerte um die Hälfte.

Die Kälte bleibt \textbf{Stärke}+\textbf{Magieniveau} Minuten bestehen.

\spellinfobox{5}{1}{0 Schritt}{Stärke+Magieniveau Minuten}{---}{1 Kampfhandlung}{Nein}


\subsection{Feuerball}

\spellheader{Schaden}{Feuer}{1 Kampfhandlung}
Der Zaubernde schleudert einen schädelgroßen lodernden Feuerball auf sein Ziel. Beim Aufprall verursacht der Zauber  \textbf{Stärke} Treffer und versetzt alles brennbare in den \textbf{Brennend} \textbf{Magieniveau} Status.

\spellinfobox{5}{1}{15 Schritt}{---}{---}{1 Kampfhandlung}{Nein}


\subsection{Ring aus Gras}

\spellheader{Heilung}{Erde}{1 Kampfhandlung}
Der Zauberne erschafft einen Ring aus Gras, welcher an dem bestimmten Ort in Sichtweite wächst und \textbf{Stufe} Meter Durchmesser hat. Das Gras sondert einen magischen Dunst ab, jeder, der den Ring durchschreitet regeneriert einmalig \textbf{Magieniveau} Wunden und ist \textit{Geschockt 1}.

Der Zauber dauert \textbf{Stärke} Kampfrunden an.

\spellinfobox{5}{1}{0 Schritt}{Stärke Runden}{Kreis}{1 Kampfhandlung}{Nein}


\subsection{Elementarschild}

\spellheader{Transmutation}{Feuer}{1 Kampfhandlung}
Der Zaubernde erschafft einen magischen Schild um seinen Körper mit seinem bevorzugten Element. Das Schild bietet dem Zaubernden im Kampf einen Schutz von \textbf{Magieniveau}×2 für \textbf{Stärke} Kampfrunden und einen zusätlichen Effekt abhängig vom Element:

Feuer: Angreifer müssen bei jedem Angriff einen W6 würfeln. bei 1-3 bekommen Sie den Zustand Brennend 1. Bei Geschossen nicht magischer Art (Pfeile, Steinschleuder, Speere, etc) besteht auch eine 50\%tige Wahrscheinlichkeit, dass die Projektile verbrennen und nicht durchdringen.

\spellinfobox{7}{2}{0 Schritt}{Stärke Runden}{---}{1 Kampfhandlung}{Nein}


\subsection{Schock}

\spellheader{Transmutation}{Energie}{1 Kampfhandlung}
Feuert einen Blitz aus reiner Energie auf einen Gegner. Der Blitz verursacht \textbf{Stärke}-3 Treffer und hat eine Durchschlagskraft von 2. Das Opfer erhält "Geschockt \textbf{Magieniveau}", auch wenn der Zauber keine Wunden verursacht.

\spellinfobox{5}{1}{10 Schritt}{---}{---}{1 Kampfhandlung}{Nein}


\subsection{Elementarer Köcher der Wandlung}

\spellheader{Beschwörung}{Natur}{1 Kampfhandlung}
Dieser magisch verzauberte Köcher erkennt den aktiven Zauber des Trägers und wandelt gewöhnliche Pfeile in elementare Projektile, die dem gewirkten Zauber entsprechen. Die Wandlung erfolgt augenblicklich beim Ziehen des Pfeils.

\spellinfobox{5}{1}{0 Schritt}{---}{---}{1 Kampfhandlung}{Nein}


\subsection{Dornensturm}

\spellheader{Schaden}{Natur}{1 Kampfhandlung}
Der Waldläufer ruft den Zorn der Natur an. Aus dem Boden schießen scharfe Dornen und peitschen durch die Luft.

Alle feindlichen Einheiten in einem Kegel von 12 Zoll Reichweite erleiden W6+2 Schaden.

Betroffene Einheiten müssen einen Rüstungswurf mit -1 Modifikator ablegen.

Getroffene Einheiten haben in ihrer nächsten Runde -2 Bewegung (die Dornen halten sie fest).

Zusatz: "Zorn des Waldes" (verstärkte Variante) Wenn der Waldläufer sich in einem bewaldeten oder naturverbundenen Gelände befindet, kann er 1 zusätzliches Mana ausgeben. In diesem Fall:

Schaden erhöht sich auf W8+3

Gegnerische Einheiten erleiden zusätzlich 1 Punkt anhaltenden Schaden zu Beginn ihrer nächsten Runde.

Gegenzauber: Ein Feind kann diesen Zauber durch eine Magieresistenzprobe (8+) negieren.

Flufftext:

„Aus jedem Blatt wird eine Klinge, aus jedem Schatten eine Falle. Die Natur vergisst nichts – und sie kämpft mit mir.“

\spellinfobox{5}{3}{12 Schritt}{---}{Kegel}{1 Kampfhandlung}{Nein}


\subsection{Aeralis Ascendere}

\spellheader{Verzauberung}{Luft}{1 Kampfhandlung}
Wenn der Zaubernde Aeralis Ascendere wirkt, antwortet die Luft selbst.

Zuerst wird die Welt still. Geräusche klingen gedämpft, als würde der Atem der Umgebung innehalten. Ein kühler Wind streicht über Haut und Kleidung, nicht zufällig, sondern bewusst – als prüfe das Element den Willen dessen, der es ruft.

Feine Strömungen sammeln sich unter den Füßen, unsichtbar und doch spürbar. Das Gewicht des Körpers beginnt zu weichen, als würde die Schwerkraft vergessen, dass sie Anspruch erhebt. Der Zaubernde erhebt sich langsam, getragen von einem Aufwind aus reinem Elementarwillen.

Im Flug fühlt sich die Bewegung nicht wie Gehen oder Springen an, sondern wie Loslassen. Gedanken lenken die Richtung, Emotionen die Geschwindigkeit. Der Wind passt sich an – mal sanft tragend, mal kraftvoll ziehend – stets im Einklang mit der inneren Ruhe des Zaubernden.

Fliegen ist kein erzwungener Akt mehr, sondern ein natürlicher Zustand. Der Körper wird Teil der Strömung, ein lebendiger Gedanke im Himmel.

Der Zauber endet lautlos. Der Wind zieht sich zurück, das Gewicht kehrt zurück – behutsam, respektvoll, als hätte die Luft den Zaubernden nur kurz ausgeliehen, um ihn an die Freiheit zu erinnern.

\spellinfobox{5}{3}{0 Schritt}{---}{Sphäre}{1 Kampfhandlung}{Nein}


\subsection{Zephyrs Griff \& Sturmschlag}

\spellheader{Kontrolle}{Luft}{1 Kampfhandlung}
Mit einer krallenden Geste verdichtet der Magier die Atmosphäre. Ein kaum wahrnehmbares Flimmern der Luft legt sich um den Gegenstand. Mit einem rissigen Geräusch, ähnlich einem Peitschenknall, bricht der Widerstand: Das Objekt schießt entweder gehorsam in die wartende Hand des Wirkenden oder wird mit der Kraft eines Orkans davonkatapultiert, wobei es beim Aufprall verheerenden Schaden anrichtet. Geübte Magier sind in der Lage, die Flugbahn des Gegenstands zu beeinflussen. Der Mindestwurf ist dabei um 2 erhöht.

\spellinfobox{10}{2}{100 Schritt}{---}{---}{1 Kampfhandlung}{Nein}


\section{Geisterbeschwörungen}

\subsection{Todesvision}

\spellheader{Weissagung}{Blut}{1 Kampfhandlung}
Das Opfer wird von einer sehr realistischen Vision seines Todes gequält, die Art des Todes kann der Zaubernde bestimmen. Die Vision umfasst das Sterben, das Verfaulen des Fleisches und den Zerfall der Knochen zu Staub. Das Opfer nimmt keinerlei körperlichen Schaden durch den Zauber, allerdings besteht die Möglichkeit, dass es durch die Todesvision traumatisiert wird. Während der Dauer des Zaubers ist das Opfer kaum zu einer sinnvollen Handlung fähig.

Das Opfer kann mit einer Willenskraftprobe versuchen, den Zauber abzubrechen. Hierbei sind so viele Erfolge notwendig wie der Zauber \textbf{Stärke} hat.

Der Zauber dauert an bis die Willenskraftprobe gelungen ist.

\begin{pquote}{Friedrich Nietzsche}
To die proudly when it is no longer possible to live proudly. Death of one's own free choice, death at the proper time, with a clear head and with joyfulness, consummated in the midst of children and witnesses: so that an actual leave-taking is possible while he who is leaving is still there.
\end{pquote}

\spellinfobox{9}{2}{0 Schritt}{Speziell Runden}{---}{1 Kampfhandlung}{Nein}


\subsection{Tiergeist rufen}

\spellheader{Beschwörung}{Natur}{Ritual}
Der Zaubernde zeichnet ein Pentagram in den Boden und konzentriert sich auf den Tiergeist.

In dem Moment in dem das Ritual beendet ist erscheint an der Stelle des Pentagrams der Tiergeist. Der Tiergeist führt \textbf{Stärke} einfache Dienste für seinen Meister aus. Die Dienste müssen einfach sein und nur eine Handlung beinhalten. z.B. einen Ritt, der beginnt mit Aufsitzen und endet mit Absteigen, oder Hilfe im Kampf gegen einen Gegner. Die Werte des Tiergeistes entsprechen denen des normalen Tieres, erhöht um \textbf{Magieniveau} Punkte.

\spellinfobox{7}{2}{0 Schritt}{---}{---}{Ritual}{Nein}


\subsection{Stimme der Toten}

\spellheader{Weissagung}{Blut}{1 Kampfhandlung}
Der Zaubernde legt einen Gegenstand des Toten vor sich, schliesst die Augen und konzentriert sich auf den Gegenstand und das Reich der Toten.

Der Zaubernde verfällt in eine leichte Trance, seine Stimme verändert sich und ähnelt der des Toten, je mehr desto persönlicher Gegenstand ist und je mehr der Zaubernde über den Toten weiß. Der Zaubernde kann \textbf{Stärke}+\textbf{Magieniveau} Fragen an den Toten stellen, die mit Ja/Nein zu beantwortet wird. Der Tote die muss Antwort auch vor seinem Tod hätte geben können.

\spellinfobox{7}{2}{0 Schritt}{---}{---}{1 Kampfhandlung}{Nein}


\subsection{Schutz der Geister}

\spellheader{Widerrufung}{Blut}{1 Kampfhandlung}
Der Zaubernde beschwört den Schutz der Geister herauf. Der Schutz von \textbf{Stärke} Personen erhöht sich für \textbf{Magieniveau} Kampfrunden um 3 Normalen Schutz, der Wert Resistenz um 3.

\spellinfobox{5}{2}{0 Schritt}{Magieniveau Runden}{---}{1 Kampfhandlung}{Nein}


\subsection{Rat der Geister}

\spellheader{Weissagung}{Blut}{1 Kampfhandlung}
Der Zaubernde kann \textbf{Stärke} Fragen an die Geisterwelt richten, die ihm, wenn die Geister ihm wohlgesonnen sind, beantwortet wird. Die Fragen müssen eine einfache ja/nein Antwort ermöglichen.

\textbf{Magieniveau} Geister erscheinen, um mögliche Fragen zu beantworten.

\spellinfobox{5}{1}{0 Schritt}{---}{---}{1 Kampfhandlung}{Nein}


\subsection{Lebende Rüstung}

\spellheader{Beschwörung}{Geist}{1 Kampfhandlung}
Der Zaubernde erschafft \textbf{Stärke} lebende Rüstungen. Die Rüstungen können einfache Verteidigungs- und Angriffsaufträge ausführen. Sie haben 4 mögliche Wunden und führen Schwerter mit einem Durchschlag von 0 und 3+\textbf{Magieniveau} Würfeln.

\spellinfobox{9}{2}{0 Schritt}{Stärke Runden}{---}{1 Kampfhandlung}{Nein}


\subsection{Körper beseelen}

\spellheader{Beschwörung}{Geist}{Ritual}
In dem Moment in dem das Ritual beendet ist, bindet der Zaubernde einen einfachen Geist in einen toten Körper, der den Körper steuert und einfache Befehle durchführt. Die Bewegungen sind langsam, und da ihm nur normale Bewegungen zur Verfügung stehen, sollte der Körper vorher besonnen aussucht werden. Eine Puppe oder Leiche kann laufen, eine Kugel kann rollen aber zum Beispiel keine Treppen steigen.

Der Zaubernde kann die direkte Kontrolle über den Körper übernehmen, als würde er in diesem stecken. Allerdings kostet dies 1 Arkana pro \textbf{Stärke}×5 Minuten, und sämtlicher Schachen, den der Körper erleidet, erleidet auch der Körper des Zaubernden.

Der Zauber endet nach \textbf{Magieniveau} Stunden.

\spellinfobox{7}{3}{0 Schritt}{Magieniveau Stunden}{---}{Ritual}{Nein}


\subsection{Grabeskälte}

\spellheader{Schaden}{Blut}{1 Kampfhandlung}
In dem Moment in dem der Zaubernde sein Ziel berührt, breitet sich von der Stelle der Berührung eine furchtbare Kälte aus, die \textbf{Magieniveau} Wunden pro Kampfrunde verursacht. Der Zaubernde kann den Zauber jederzeit abbrechen, wird er nicht abgebrochen, endet er wenn das Ziel vollständig ausgekühlt und tot ist.

Das Opfer wirft zu Beginn jeder Kampfrunde vor dem Schaden einen Wurf auf Resistenz. Gelingt der Wurf mit mindestens \textbf{Stärke} Erfolgen, so ist der Zauber beendet und es entsteht kein Schaden mehr.

\spellinfobox{9}{2}{0 Schritt}{---}{---}{1 Kampfhandlung}{Nein}


\subsection{Geist austreiben}

\spellheader{Widerrufung}{Geist}{1 Kampfhandlung}
Der Zaubernde wählt bis zu \textbf{Stärke} Geister oder von Geistern beherrschte Wesen. Die Geister werden verbannt und verlassen ihr irdisches Dasein.

\spellinfobox{5}{3}{0 Schritt}{---}{---}{1 Kampfhandlung}{Nein}


\subsection{Gegenstand beseelen}

\spellheader{Beschwörung}{Geist}{Ritual}
In dem Moment in dem das Ritual beendet ist, bindet der Zaubernde einen einfachen Geist in den gewählten Gegenstand, der eine einfache Aktion an dem Gegenstand durchführt.

Im Gegensatz zu dem Ritual Beseelte Waffe ist der Grundgedanke dieses Rituals eher friedlich, so dass die verbreitesten Einsatzbereiche diesen Rituals dazu dienen um in Kristallkugeln leuchtenden Nebel wirbeln zu lassen, oder Kerzenleuchter die Kerzen entzünden zu lassen, sobald der Raum betreten wird.

Bei dem Ritual muss der Zaubernde festlegen, wer und wie der Auslöser ist. Für manche Sachen, wie die genannten Kristallkugeln, macht es Sinn, wenn dies nur eine bestimmte Berührung des Nutzers ist, bei dem angesprochenem Kerzenleuchter eher jeder, der in eine bestimmte Reicheite um den Kerzenleuchter kommt.

Der Gegenstand bleibt bis zu \textbf{Stärke}+\textbf{Magieniveau} Stunden beseelt.

\spellinfobox{5}{5}{0 Schritt}{Stärke+Magieniveau Stunden}{---}{Ritual}{Nein}


\subsection{Einfacher Dienstgeist}

\spellheader{Beschwörung}{Natur}{Ritual}
Der Zaubernde reibt seine Hand mit Knochenstaub ein, zeichnet ein Pentagram in die Luft und konzentriert sich auf den Geist.

In dem Moment in dem das Ritual beendet ist erscheint an der Stelle des Pentagrams der Dienstgeist. Der Dienstgeist führt einen einfachen Dienst für seinen Meister aus, der Dienst kann aus max. \textbf{Stärke} Dingen bestehen. Der Dienstgeist ist nicht in der Lage einem Lebewesen oder Gegenstand direkt Schaden zuzufügen. Beispiele für Dienste sind das Überbringen von sehr kurzen Nachrichten (max. \textbf{Stärke} Worte an eine Person, oder 1 Wort an insgesamt \textbf{Stärke} Personen), die Benachrichtigung des Zaubernden wenn eins von \textbf{Stärke} bestimmten Ereignissen eintritt oder gar das Sammeln von \textbf{Stärke} Äpfeln.

Magieniveau 4+: Der Geist kann auch Personen schaden zufügen.

\spellinfobox{7}{1}{0 Schritt}{---}{---}{Ritual}{Nein}


\subsection{Besessenheit}

\spellheader{Kontrolle}{Geist}{Ritual}
Der Zaubernde zeichnet dem Ziel ein Pentagram auf die Stirn und konzentriert sich auf den Geist und das Ziel.

Bei diesem Ritual lässt der Zaubernde einen Geist in den Körper des Ziels fahren. Der Geist kann ein Dienstgeist, ein freier Geist oder der Geist des Zaubernden sein. Ihm letzten Fall sackt der Körper des Zaubernden zusammen, reagiert nicht, atmet langsam und starrt in die Leere, wenn seine Augen geöffnet werden. Das Ziel befindet sich (\textbf{Stärke}+\textbf{Magieniveau})×10 Minuten unter der Kontrolle des in es gefahrenen Geistes, der den gesamten Körper steuern kann (laufen, schlagen, kämpfen, sprechen, usw.).

Sollte der Körper des Ziels sterben während der Bessesenheit, so verschwindet der Geist schlagartig aus dem Körper und fährt in seine Sphäre zurück. Sollte der Körper des Ziels sterben während der Zaubernde in ihm weilt, so kehrt der Geist des Zaubernden in seinen ursprünglichen Körper zurück, und der Zaubernde ist für 3W6 Minuten bewusstlos.

\spellinfobox{11}{3}{0 Schritt}{(Stärke+Magieniveau)×10 Minuten}{---}{Ritual}{Nein}


\subsection{Beseelte Waffe}

\spellheader{Beschwörung}{Natur}{Ritual}
Der Zaubernde formt einen Beschwörungskreis um die zu beseelende Waffe.

In dem Moment in dem das Ritual beendet ist, bindet der Zaubernde einen Geist an die beseelte Waffe. Dieser Geist fügt dem Opfer, neben den waffenüblichen Wunden, zusätzlich \textbf{Stärke} Wunden zu. Die Waffe gilt als magische Waffe, kann jedoch nur von dem Zaubernden verwendet werden. Jeder andere der versucht die Waffe zu bedienen, wird selbst von dem Geist attackiert. Das Erscheinungsbild des Geistes kann vom Zaubernden frei bestimmt werden.

Mit der Waffe kann \textbf{Magieniveau}×5 mal angegriffen werden, bevor der Geist die Waffe verlässt.

\spellinfobox{15}{5}{0 Schritt}{Magieniveau×5 Aktionen}{---}{Ritual}{Nein}


\subsection{Geistkörper}

\spellheader{Transmutation}{Geist}{1 Kampfhandlung}
Der Zaubernde schließt die Augen, denkt den Vers "Mein Körper, ein Geist" und öffnet die Augen danach wieder.

Der Zaubernde ist in der Lage, sämtliche Aktionen eines Geistes durchzuführen, z.B. sehen, Dinge berühren, zaubern sofern es keine Zutaten benötigt oder diese in seiner Reichweite liegen, fliegen, sich durch unbelebte Gegenstände bewegen usw.

Der Zauber kann jederzeit fallengelassen werden, wodurch der Zaubernde seine normale körperliche Gestalt annimmt, allerdings ohne Kleidung, da diese nicht mit verwandelt wird. Dem Zaubernden stark vertraute Gegenstände sollen wohl auch in Geisterform mitgeführt werden. Da Geister gegen nichtmagische Waffen und Attacken immun sind, ist dies auch der Zaubernde.

Der Zaubernde kann in Geistgestalt auch gebannt werden. Wenn er ausgetrieben wird, landet der Zaubernde in seinem Körper, bewusstlos für 2W6 Minuten in seinem Körper, an der Stelle wo er den Zauber begonnen hat.

Der Zauber hält für \textbf{Stärke}+\textbf{Magieniveau} Minuten an.

\spellinfobox{4}{1}{0 Schritt}{Stärke+Magieniveau Minuten}{---}{1 Kampfhandlung}{Nein}


\section{Hermetik}

\subsection{Ungesehen, unbemerkt}

\spellheader{Kontrolle}{Geist}{1 Kampfhandlung}
Der Zaubernde erscheint seiner Umgebung im wahrsten Sinne des Wortes als ein Nichts. Ihm wird keinerlei Beachtung geschenkt, Menschen rempeln ihn auf der Straße an, kümmern sich jedoch nicht darum. Selbst wenn er jemanden anspricht, wird er ignoriert. Schafft er es, die Aufmerksamkeit von jemandem auf sich zu lenken, so vergißt dieser ihn sofort wieder, sobald er sich ihm entzieht.

Bei diesem Zauber handelt es sich nicht um eine Verwandlung des Zaubernden, sondern um eine Massenhypnose. Das hat zur Folge, daß der Zaubernde auch durch Hellsicht-Zauber wie Leben erkennen nicht zu entdecken ist. Seine Aura ist zwar genauso sichtbar wie eh und je, ihr wird jedoch keine Beachtung geschenkt.

Jemandem, der gezielt nach dem Zaubernden sucht, steht eine Wahrnehmungsprobe zu, um ihn dennoch zu entdecken. Gelingt diese mit \textbf{Stärke} Erfolgen, so fällt der Hypnose-Effekt von dem Suchenden ab und er kann den Zaubernden wieder normal wahrnehmen.

Dem Zaubernden muß nach dem Zauber eine Willenskraftprobe gegen den Mindestwurf \textbf{Magieniveau}×2 gelingen, um nicht in eine tiefe Depression zu verfallen.

\spellinfobox{7}{2}{0 Schritt}{Stärke Stunden}{---}{1 Kampfhandlung}{Ja}


\subsection{Schleier des Vergessens}

\spellheader{Kontrolle}{Geist}{1 Kampfhandlung}
Dieser Zauber erlaubt es dem Zaubernden, die Erinnerungen seines Opfers an ein Ereignis der letzten \textbf{Stärke} Stunden zu manipulieren. Das Opfer vergißt für \textbf{Magieniveau}×10 Minuten eine vom Zaubernden festgelegte, mit dem betreffenden Ereignis in Verbindung stehende Aufgabe auszuführen (zum Beispiel, die Wachen zu alarmieren). Wird es durch irgendetwas oder -jemanden an die Aufgabe erinnert, fällt der Zauber sofort von ihm ab.

Das Opfer darf einen Wurf auf Willenskraft machen. Erreicht es dabei Erfolge entsprechend der \textbf{Stärke} des Zaubers so wird es nicht manipuliert.

\spellinfobox{5}{2}{0 Schritt}{Magieniveau×10 Minuten}{---}{1 Kampfhandlung}{Nein}


\subsection{Öffnen}

\spellheader{Transmutation}{Energie}{1 Kampfhandlung}
Der Zaubernde ist in der Lage verschlossene, nicht magische Objekte wie Türen, Truhen oder andere Schlösser zu öffnen.

Zusätzlich zu normal verschlossenen Schlössern ist es dem Zaubernden möglich, magisch verschlossene Schlösser mit einem Siegel von der Stufe der \textbf{Stärke des Zaubers} zu öffnen.

Magieniveau 5+: Der Zauber öffnet alle Schlösser.

\spellinfobox{3}{1}{0 Schritt}{---}{---}{1 Kampfhandlung}{Nein}


\subsection{Magieanalyse}

\spellheader{Weissagung}{Arkana}{1 Kampfhandlung}
Der Zaubernde ist in der Lage eine Analyse eines gewirkten oder gerade im wirken begriffenen Zaubers durchzuführen. Der Zaubernde erkennt die Schule der Magie, Wesen des Zaubers (ob nun Heilung, Schaden, Art des Elements, Dauer usw.) sowie eine grobe Einschätzung der Stärke des Zaubers.

\spellinfobox{5}{1}{40 Schritt}{---}{---}{1 Kampfhandlung}{Nein}


\subsection{Geist des Weines}

\spellheader{Kontrolle}{Geist}{1 Kampfhandlung}
Der Zaubernde deutet an, ein Glas Wein zu trinken. Dabei murmelt er den Namen des Zaubers.

Bis zu \textbf{Magieniveau} Opfer des Zaubers erfahren augenblicklich einen Rausch der Trunkenheit, der sie torkeln lässt und jede normale Handlung erschwert. der Mindestwurf für alle Würfe ist um 2 erhöht. Die Wirkung des Zaubers hält maximal \textbf{Stärke}×10 Minuten an.

Jedes Opfer des Zaubers darf einen Wurf auf Willenskraft machen. Erreicht der Wurf Erfolge entsprechend der \textbf{Stärke} des Zaubers, so widersteht es dem Zauber.

\spellinfobox{7}{3}{0 Schritt}{Stärke×10 Minuten}{---}{1 Kampfhandlung}{Nein}


\subsection{Extreme Leistung}

\spellheader{Kontrolle}{Blut}{1 Kampfhandlung}
Der Zaubernde kann für eine schwierige Aufgabe ein Persona Attribut, eine Kampffertigkeit (Schießen, Nahkampf oder Werfen) oder 'Ausweichen' kurzfristig extrem steigern. Der gewählte Wert steigt für einen Zeitraum von \textbf{Stärke} Minuten um \textbf{Magieniveau}×2 Punkte.

\spellinfobox{7}{2}{0 Schritt}{Stärke Minuten}{---}{1 Kampfhandlung}{Nein}


\subsection{Erstarren}

\spellheader{Transmutation}{Wasser}{1 Kampfhandlung}
Das Opfer erstarrt an Ort und Stelle für \textbf{Stärke} Aktionen. Es ist bei vollem Bewusstsein und alle Sinne funktionieren normal. Körperliche Handlungen oder Angriffe sind jedoch nicht möglich.

Zu Beginn jeder Aktion des Opfers wirft dieses auf Willenskraft, wobei der Mindestwurf um das \textbf{Magieniveau} erhöht wird. Ist der Wurf erfolgreich, wird die Starre aufgehoben und die Aktion steht dem Opfer zur Verfügung. Die Erschwernis des Mindestwurfes sinkt nach jedem Wurf um 1.

\spellinfobox{5}{2}{0 Schritt}{Stärke Aktionen}{---}{1 Kampfhandlung}{Nein}


\subsection{Dilatio}

\spellheader{Transmutation}{Luft}{1 Kampfhandlung}
Der Zaubernde erschafft ein kurzzeitiges spontanes Portal unter sich, in das er sofort hinein gesaugt wird. Ein weiteres Portal erschafft er an einem Ort, welcher sich maximal \textbf{Stärke}×10 Schritt von seiner aktuellen Position entfernt befindet.

Ohne jede Verzögerung erscheint er am gewünschten Zielort.

\spellinfobox{7}{2}{0 Schritt}{---}{---}{1 Kampfhandlung}{Nein}


\subsection{Blutrausch}

\spellheader{Kontrolle}{Blut}{1 Kampfhandlung}
Der Verzauberte fällt in einen unkontrollierbaren Blutrausch. Seine geistigen Fähigkeiten sind soweit verkümmert, dass er gerade noch Freund von Feind unterscheiden kann. Seine Kampfwerte (Schießen, Nahkampf und Werfen) steigen jeweils um \textbf{Stärke} Punkte.

Bildung, Logik und Geschick sinken um \textbf{Magieniveau} Punkte. Der Verzauberte empfindet keinen Schmerz und keine Erschöpfung, lediglich eine unbeherrschbare Kampfeslust. Nachdem der Zauber von ihm gefallen ist, bricht er bewusstlos zusammen.

\spellinfobox{7}{1}{0 Schritt}{Stärke Runden}{---}{1 Kampfhandlung}{Nein}


\subsection{Arrest}

\spellheader{Verzauberung}{Arkana}{1 Kampfhandlung}
Das Opfer der Verzauberung ist für \textbf{Stärke}+\textbf{Magieniveau} Kampfrunden an dem Ort gefangen, an dem es sich befindet. Es ist in der Lage normal zu handeln und auch anzugreifen, kann sich jedoch nicht vom Fleck bewegen.

\spellinfobox{5}{2}{0 Schritt}{Stärke Runden}{---}{1 Kampfhandlung}{Nein}


\subsection{Unsichtbarkeit}

\spellheader{Transmutation}{Licht}{1 Kampfhandlung}
Der Zaubernde hüllt sich selbst oder ein anderes Wesen in einen arkanen Mantel aus reflektierendem Licht, so dass seine Gestalt nicht zu erkennen ist.

Für die Dauer des Zaubers erhält das Ziel \textbf{Stärke} Punkte auf die Fertigkeit Heimlichkeit.

\spellinfobox{5}{1}{1 Schritt}{Magieniveau Minuten}{---}{1 Kampfhandlung}{Nein}


\subsection{Magie Absorbieren}

\spellheader{Widerrufung}{Arkana}{1 Kampfhandlung}
Der Zaubernde kann Arkana aus einem magischen Gegenstand oder einem Magischen Lebewesen absorbieren. Es werden \textbf{Stärke}+\textbf{Magieniveau} Arkana vom Ziel auf den Zaubernden übertragen.

Ein magisches Lebewesen wirft auf seine Willenskraft. Für jeden Erfolg bei diesem Wurf wird die Menge an übertragenem Arkana um 1 verringert.

\spellinfobox{5}{1}{15 Schritt}{---}{---}{1 Kampfhandlung}{Nein}


\subsection{Aevum}

\spellheader{Transmutation}{Geist}{Ritual}
Das Ritual bedarf einer ausführlichen Vorbereitung. So ist sowohl die genaueste Vorbereitung des Hermetikers als auch die der Zielperson erforderlich.

Der Hermetiker aktiviert den Zauber bereits zu Beginn des Rituals. Über die ganze Zeit muss er diesen Zauber aufrecht erhalten, was eine erhebliche Menge Magie verschlingt. An jedem einzelnen Tag muss er zwei Stunden der Meditation aufbringen, wobei er an dem Gemälde arbeitet. Das Gemälde muss ausschließlich vom Hermetiker erschaffen werden.

Die Zielperson benötigt keine weitere Vorbereitung als die Erkenntnis fleischlos zu werden. Zu diesem Zwecke sollte er sich selbst mit allen möglichen Arten der Verbrennung, Vergiftung und Ähnlichem martern, um den Abschluss des Rituals zu vereinfachen.

Zum Zeitpunkt der Ausführung muss der Hermetiker lediglich die Sphäre mit seiner Konzentration aufrecht erhalten und die Arme steuern. Szenerie sowie Atmosphaere sind wie bei den meisten hermetischen Handlungen unerheblich.

Über die Zeit der Erschaffung des Bildes bindet der Hermetiker den Geist und die Seele der Zielperson an dieses. Zur Zeit der Ausführung entsteht unweit des Gemäldes eine Sphäre aus reiner Magie, welche in einem halbdurchsichtigen matten Weiss zumeist in der Luft schwebt. Diese Sphäre bildet Arme aus, welche Schläuchen gleich über die Köpfe der bereitzustellenden Opfer gleiten. Unter einem fortwährenden Brummen sammelt die Sphäre die Innereien der Opfer, um sie danach im Bild zu verdichten.

Der Hermetiker erschafft so eine feste Bindung des Geistes und der Seele der Zielperson an das Gemälde. Der Geist im Gemälde ist jederzeit in der Lage innerhalb einer Sekunde in eine beliebige Person in direkter Nähe des Gemäldes einzufahren. In dieser Person lebt die Zielperson dann bis zu deren Tode oder einer Austreibung fort, um danach erneut in das Bild einzukehren. Sie beherrscht die Person völlig, empfindet deren Gefühle und lebt vollkommen in ihr. Dieser Fortgang des Beherrschens und Zurückfahrens nimmt erst ein Ende sobald das Bild zerstört ist oder die Zielperson vier Mal in das Bild zurückgekehrt ist. Danach verbleibt er im Bild.

Das erschaffene Bild ist selbst nahezu unzerstörbar. Es gilt als Artefakt der Stufe des Hermetikers und bedarf des selben Aufwands zur Vernichtung wie jedes andere Artefakt dieser Stufe. Einzig ein immenses Magieeinwirken oder göttliches Wirken können das Objekt zerstören.

Die Zielperson kann \textbf{Stärke}+\textbf{Magieniveau} Mal in das Bild einfahren, bevor sie darin gefangen ist.

\spellinfobox{80}{1000}{0 Schritt}{---}{---}{Ritual}{Nein}


\subsection{Expolitio}

\spellheader{Transmutation}{Wasser}{1 Kampfhandlung}
Der Zaubernde vermag es, eine Fläche von \textit{Stärke des Zaubers} Schritt im Quadrat von jeglicher Verschmutzung zu reinigen. Die Reinigung findet sofort statt und entfernt Schmutz und Gerüche.

\spellinfobox{3}{1}{0 Schritt}{---}{---}{1 Kampfhandlung}{Nein}


\subsection{Arkaner Schutz}

\spellheader{Transmutation}{Arkana}{1 Kampfhandlung}
Für die Dauer des Zaubers kann der Zaubernde Arkana ausgeben, als sei es Schutz. Jedes im Sinne der Schutz Regel ausgegebene Arkana zählt als ein "Schutz vor kritischen Treffern".

\spellinfobox{5}{1}{0 Schritt}{Magieniveau Runden}{---}{1 Kampfhandlung}{Nein}


\section{Mystizismus}

\subsection{Translokation}

\spellheader{Kontrolle}{Geist}{Ritual}
Der Zaubernde beginnt nach der Vorbereitungszeit von 30 Sekunden langsam zu verschwimmen und undeutlicher zu werden. Dieser Prozess dauert weitere 30 Sekunden, in denen der Zaubernde jedoch nicht mehr verletzbar ist. Versucht jemand, den Zaubernden in dieser Zeit zu berühren, so durchdringt seine Hand eine kalte, dichte Masse.

Der Zaubernde kann sich maximal \textbf{Stärke} Kilometer weit trannsportieren lassen. Sein Körper taucht am Ziel auf, ohne wirklich den Weg zurückzulegen. Am Zielort dauert es widerum 30 Sekunden, bis der Körper ganz erschienen ist.

\spellinfobox{2}{3}{0 Schritt}{---}{---}{Ritual}{Nein}


\subsection{Schutzaura}

\spellheader{Beschwörung}{Arkana}{1 Kampfhandlung}
Der Zaubernde erschafft eine magische Schutzaura um sich herum. Die Aura strahlt von ihm ab und ist durch magische Handlungen zu erkennen. Sie ist für alle Formen der Magie undurchdringlich, worunter auch magische Angriffe oder Verwandlungen fallen. Die Schutzaura kann \textbf{Stärke}×2 Wunden absorbieren, bevor sie in sich zusammenfällt. Der Zauber kann maximal \textbf{Magieniveau} Kampfrunden aufrechterhalten werden.

\spellinfobox{5}{1}{0 Schritt}{Magieniveau Runden}{---}{1 Kampfhandlung}{Nein}


\subsection{Falsches Gefühl}

\spellheader{Kontrolle}{Geist}{1 Kampfhandlung}
Das Opfer empfindet ein Gefühl, welches vom Zaubernden gewählt wurde. Dabei handelt es sich um ein konkretes Gefühl, wie "Stolz auf etwas". Das Gefühl hält \textbf{Stärke} Minuten an.

\spellinfobox{7}{1}{0 Schritt}{Stärke Minuten}{---}{1 Kampfhandlung}{Nein}


\subsection{Schnelligkeit}

\spellheader{Verzauberung}{Luft}{1 Kampfhandlung}
Der Zaubernde beschleunigt sich selbst. Seine Schnelligkeit sowie die Reichweiten für \textit{Rennen} und \textit{Gehen} werden für \textbf{Magieniveau} W6 Runden um \textbf{Stärke} erhöht.

\spellinfobox{3}{1}{0 Schritt}{Magieniveau W6 Runden}{---}{1 Kampfhandlung}{Nein}


\section{Nekrologie}

\subsection{Zombie erwecken}

\spellheader{Beschwörung}{Arkana}{1 Kampfhandlung}
Der Zaubernde erweckt bis zu \textbf{Magieniveau} Leichen in einem Umkreis von 10 Metern. Die Zombies folgen seinem Befehl und bleiben bis zu \textbf{Stärke} Minuten lebendig.

\spellinfobox{11}{3}{0 Schritt}{Stärke Minuten}{---}{1 Kampfhandlung}{Nein}


\subsection{Weg der Knochen}

\spellheader{Transmutation}{Blut}{1 Kampfhandlung}
Der Zaubernde zerfällt zu feinem Staub und kann in dieser Form mit 10 Kilometern / Stunde reisen. Maximal kann er sich in dieser Form \textbf{Stärke} Stunden halten.

Der Zaubernde hat in dieser Form \textbf{Magieniveau} Wunden, kann allerdings nur von Dingen verwundet werden, die einem Haufen Knochen schaden machen.

\spellinfobox{3}{2}{0 Schritt}{Stärke Stunden}{---}{1 Kampfhandlung}{Nein}


\subsection{Toter Wegweiser}

\spellheader{Weissagung}{Blut}{1 Kampfhandlung}
Ist in einem Umkreis von \textbf{Stärke}×10 Metern eine Leiche (auch Tierleichen gelten, sofern sie mindestens die Grösse einer Maus haben) im Boden vergraben, so ist es dem Zaubernden gestattet \textbf{Magieniveau} Fragen nach einer Richtung zu stellen: "In welcher Richtung liegt Meridian?" "Der letzte Reiter der diesen Boden passiert, wohin ist er geritten?"

\spellinfobox{5}{1}{10 Schritt}{---}{---}{1 Kampfhandlung}{Nein}


\subsection{Toter Blick}

\spellheader{Weissagung}{Blut}{1 Kampfhandlung}
Der Zaubernde legt Daumen und Zeigefinger in die Augen des Toten und schließt seine Augen.

Der Zaubernde blickt rückwärts vom Zeitpunkt des Todes in die Vergangenheit des Toten. Dabei sieht er alles, das der Tote aus seinen Augen gesehen hat. Die Bilder wirken verschwommener, je weiter der Blick in die Vergangenheit des Toten geht. Der Zaubernde sieht maximal die letzten \textbf{Stärke} Tage vor dem Tod. Je höher das \textbf{Magieniveau}, desto klarer sind die Eindrücke.

\spellinfobox{7}{1}{0 Schritt}{---}{---}{1 Kampfhandlung}{Nein}


\subsection{Totenwesen}

\spellheader{Transmutation}{Blut}{1 Kampfhandlung}
Der Zaubernde erschafft ein untotes Mischwesen aus verschiedenen Skeletten. Das Wesen ist dazu in der Lage, unkontrollierte Handlungen auszuführen. Es kann Waffen mit einem Fertigkeitswert von \textbf{Stärke}×2 führen. Das Wesen erhält alle Fähigkeit der ehemaligen Lebewesen.

Das Wesen bleibt \textbf{Magieniveau}×2 Kampfrunden bestehen.

\spellinfobox{7}{2}{0 Schritt}{Magieniveau×2 Runden}{---}{1 Kampfhandlung}{Nein}


\subsection{Totenschwert}

\spellheader{Beschwörung}{Blut}{1 Kampfhandlung}
Der Zaubernde erschafft ein Totenschwert aus dem Nichts. Das Schwert hat einen Durchschlag von 1 und ein Schadenspotential von \textbf{Magieniveau} Würfeln im Nahkampf.

Jede verursachte Wunde absorbiert das Schwert. Für jede absorbierte Wunde erhält das Schwert einen Würfel Schadenspotential. Erreichen die absorbierten Wunden \textbf{Stärke}+1, so zerfällt das Schwert und der Zauber endet.

\spellinfobox{5}{2}{0 Schritt}{---}{---}{1 Kampfhandlung}{Nein}


\subsection{Totenschrei}

\spellheader{Illusion}{Blut}{1 Kampfhandlung}
Dem Zaubernden entfährt ein grausamer Totenschrei, der alle Anwesenden im Umkreis von 10 Metern die nicht ihre Ohren zugehalten haben, \textbf{Magieniveau} Kampfrunden entsetzt handlungsunfähig werden lässt.

Jedes Opfer wirft einen Wurf auf seine Willenskraft. Erreicht es hierbei nicht Erfolge entsprechend der \textbf{Stärke des Zaubers}, so flieht das Opfer in Panik.

\spellinfobox{6}{3}{0 Schritt}{---}{Kreis}{1 Kampfhandlung}{Nein}


\subsection{Tanz des Todes}

\spellheader{Beschwörung}{Blut}{1 Kampfhandlung}
Im Umkreis von \textbf{Stärke}×3 Metern um den Zaubernden erheben sich alle Toten aus der Erde, und stehen unter der Kontrolle des Zaubernden. Die Toten sind dazu in der Lage, unkontrollierte Handlungen auszuführen und bleiben für \textbf{Magieniveau} Stunden lebendig.

\spellinfobox{9}{3}{0 Schritt}{Stärke Stunden}{---}{1 Kampfhandlung}{Nein}


\subsection{Skelette beschwören}

\spellheader{Beschwörung}{Blut}{1 Kampfhandlung}
\textbf{Stärke} Skelette erheben sich aus dem Boden, bewaffnet mit Knochenschwertern (Durchschlag 0). Sie können mit der Fertigkeit des Zaubernden kämpfen. Sie haben \textbf{Magieniveau} mögliche Wunden und bleiben für \textbf{Stärke}×3 Kampfrunden.

\spellinfobox{9}{3}{0 Schritt}{Stärke×3 Runden}{---}{1 Kampfhandlung}{Nein}


\subsection{Pein der Knochen}

\spellheader{Kontrolle}{Blut}{1 Kampfhandlung}
Das Opfer erleidet unglaubliche Schmerzen, glaubt ihm würden seine Knochen bersten.

Misslingt dem Opfer ein Willenskraft-Wurf, so lässt der Effekt das Opfer für \textbf{Stärke} Kampfrunden handlungsunfähig am Boden zuammenbrechen. Der Mindestwurf für diesen Wurf ist um \textbf{Magieniveau} erhöht.

\spellinfobox{5}{1}{0 Schritt}{Stärke Runden}{---}{1 Kampfhandlung}{Nein}


\subsection{Odem der Vergänglichkeit}

\spellheader{Transmutation}{Blut}{1 Kampfhandlung}
Alle lebenden Pflanzen im Umkreis von \textbf{Stärke}×3 Metern gehen unter dem Atem des Zaubernden ein und verwittern zu einer toten, schwarzen Masse. Dies betrifft auch magisch geschaffene Pflanzen.

\spellinfobox{3}{1}{0 Schritt}{---}{---}{1 Kampfhandlung}{Nein}


\subsection{Knochenpeitsche}

\spellheader{Beschwörung}{Blut}{1 Kampfhandlung}
Aus dem Arm des Zaubernden formt sich eine ca. 3 Meter lange Knochenpeitsche, die von dem Zaubernden als Waffe geführt werden kann. Die Reichweite der Peitsche beträgt 3 Meter, sie hat einen Durchschlag von 0 und ein Schadenspotential von \textbf{Magieniveau} Würfeln.

Die Peitsche bleibt bis zu \textbf{Stärke}×3 Kampfrunden bestehen.

\spellinfobox{5}{2}{0 Schritt}{Stärke×3 Runden}{---}{1 Kampfhandlung}{Nein}


\subsection{Der Zorn der Knochen}

\spellheader{Schaden}{Blut}{1 Kampfhandlung}
Der Zaubernde wirft den Gegnern Knochensplitter entgegen und beschwört den Zorn der Knochen. Die Splitter verursachen \textbf{Stärke}+1 Treffer bei allen Gegnern im Kegel, die gleichmäßig auf \textbf{Magieniveau} Gegner aufgeteilt werden.

\spellinfobox{5}{1}{20 Schritt}{---}{Kegel}{1 Kampfhandlung}{Nein}


\subsection{Geheimnisse der Toten}

\spellheader{Weissagung}{Geist}{Ritual}
Der Zaubernde kniet über einer Leiche oder einem Skelett und beginnt diese langsam zu zerlegen.

Der Zaubernde erfährt mit jedem Stück welches er schneidet ein Teil des Wissens des Toten. Nach Abschluss des Rituals hat er einen Eindruck aller Erinnerungen des Toten. Der Prozess ist für die Seele des Toten eine Tortur im Totenreich. Der Nekrologe erhält für \textbf{Stärke} Stunden +\textbf{Magieniveau} in dem Attribut, welches dem Höchsten des Toten entspricht.

\spellinfobox{4}{1}{0 Schritt}{Stärke Stunden}{---}{Ritual}{Nein}


\subsection{Knochenschild}

\spellheader{Transmutation}{Geist}{1 Kampfhandlung}
Der Nekrologe wirft einen Knochensplitter zu Boden und ruft "Schütze mich!"

Aus dem Boden schiessen Knochenstücke hervor und umgeben den Nekrologen. Praktisch kommen sie einem Schutz von \textbf{Magieniveau}+2 auf den ganzen Körper gleich. Der Knochenschild entspricht einer Belastung von 2.

Der Knochenschild bleibt \textbf{Stärke} Kampfrunden bestehen.

\spellinfobox{5}{2}{0 Schritt}{Stärke Runden}{---}{1 Kampfhandlung}{Nein}


\subsection{Knochenwuchs}

\spellheader{Beschwörung}{Geist}{1 Kampfhandlung}
Der Zaubernde hat sich in das Zentrum des zu verzaubernden Geländes zu knien, und bohrt einen Knochensplitter in den Boden.

Der Zaubernde erzeugt ein unnatürliches Gewächs aus Knochen, dass innerhalb von 25 Minuten bis zur kompletten Grösse wächst. Innerhalb eines Radius von max. \textbf{Stärke} Schritt entsteht ein grosses Gewächs aus Knochen. Der Nekrologe kann mit Hilfe des Knochenwuchs Mauern erklimmen, oder es Spalten in Türen aufbrechen lassen. Der Knochenwuchs bleibt \textbf{Magieniveau} Stunden bestehen, danach zerfällt er zu Staub. Der Zaubernde gibt dem Wuchs lediglich eine ungefähre Richtung an, kontrollieren kann er ihn nicht.

\spellinfobox{4}{1}{0 Schritt}{Magieniveau Stunden}{---}{1 Kampfhandlung}{Nein}


\section{Sanguine Magie}

\subsection{Transfusion}

\spellheader{Schaden}{Blut}{1 Kampfhandlung}
Der Zaubernde zapft die Lebensenergie eines Wesens an, hierzu stellt er mithilfe eines Blutkristalles eine Verbindung zu seinem Opfer her. Das Opfer muss hierzu eine Wunde haben, über die der Begabte die Verbindung herstellen kann. Er spinnt einen Faden aus Blut von seinem Blutkristall ausgehen zu der Wunde des Opfers.

Das Ziel erleided \textbf{Stärke} Wunden und der Zaubernde wird um die selbe Anzahl + \textbf{Magieniveau} geheilt.

\spellinfobox{7}{3}{2 Schritt}{---}{---}{1 Kampfhandlung}{Nein}


\subsection{Ritus des Lebens}

\spellheader{Heilung}{Blut}{1 Kampfhandlung}
Der Zaubernde nimmt \textbf{Stärke} Wunden hin. Das Ziel wird um die doppelte Anzahl + \textbf{Magieniveau} geheilt.

\spellinfobox{3}{2}{0 Schritt}{---}{---}{1 Kampfhandlung}{Nein}


\subsection{Ritus der Reinigung}

\spellheader{Heilung}{Blut}{Ritual}
Der Zaubernde begibt sich in einen meditativen Zustand und konzentriert sich auf schädliche Stoffe in seinem Blut, wie etwa Krankheiten und Gifte. Dieser sammelt er an einer bestimmten Stelle seines Körpers und lässt sie dann über einen Schnitt austreten.

Für jede Krankheit und jedes Gift das er kurieren möchte nimmt der Zaubernde eine Wunde hin. Diese Anzahl ist um die \textbf{Stärke} des Zaubers verringert.

\spellinfobox{5}{1}{0 Schritt}{---}{---}{Ritual}{Nein}


\subsection{Blutpfeil}

\spellheader{Schaden}{Blut}{1 Kampfhandlung}
Der Zaubernde beschwört aus seinem Blut ein pfeilartiges Geschoss, welches er dem Gegner mit hoher Geschwindigkeit entgegenschleudert.

Der Zaubernde nimmt \textbf{Magieniveau} Wunden hin. Das Opfer erleidet \textbf{Stärke} Wunden.

\spellinfobox{5}{1}{0 Schritt}{---}{---}{1 Kampfhandlung}{Nein}


\subsection{Blutbann}

\spellheader{Kontrolle}{Blut}{1 Kampfhandlung}
Der Zaubernde beeinflusst die Blutzirkulation seines Opfers und kann so Taubheit und Lähmung von bestimmten Körperregionen hervorrufen. Fertigkeiten, die das betreffende Körperteil erfordern, sind um \textbf{Magieniveau} reduziert.

Nach \textbf{Stärke} Kampfrunden muss der Zaubernde jeweils ein Arkana oder eine Wunde aufwenden um den Effekt aufrecht zu erhalten.

\spellinfobox{7}{1}{0 Schritt}{Stärke Runden}{---}{1 Kampfhandlung}{Nein}


\subsection{Ritus des Blutkristalls}

\spellheader{Transmutation}{Blut}{Ritual}
Der Zaubernde begibt sich in einen meditativen Zustand und konzentriert sich auf den Fluss der Magie und des Blutes. Nach einem Tag öffnet der Zaubernde seine Adern und lässt fast sein gesamtes Blut in eine Tonschale fliessen.

Nach der Zugabe von Rubinstaub reduziert er dieses auf die absolute Essenz und formt daraus einen Kristallsplitter. Diesen setzt er sich dann in einem Schnitt an seinem Körper ein. Über diesen Kristall wirkt der Zaubernde dann seine Sanguine Magie, ohne sich jedesmal eine separate Wunde zufügen zu müssen.

\spellinfobox{8}{4}{0 Schritt}{---}{---}{Ritual}{Nein}


\section{Schamanismus}

\subsection{Einfache Heilung}

\spellheader{Heilung}{Natur}{1 Kampfhandlung}
Der Zaubernde heilt eine Person oder ein Tier für \textbf{Stärke}+\textbf{Magieniveau} Wunden.

\spellinfobox{5}{1}{5 Schritt}{---}{---}{1 Kampfhandlung}{Nein}


\subsection{Wettervorhersage}

\spellheader{Weissagung}{Natur}{Ritual}
Der Zaubernde ist in der Lage, das Wetter für die nächsten \textbf{Magieniveau} Tage vorherzusagen. Ist die \textbf{Stärke} des Zaubers über 3, kann er die Wetteränderungen immer zeitgenau vorhersagen, andernfalls ist ihm nur bekannt, wie sich das Wetter entwickeln wird.

\spellinfobox{2}{1}{0 Schritt}{---}{---}{Ritual}{Nein}


\subsection{Vogelruf}

\spellheader{Beschwörung}{Natur}{1 Kampfhandlung}
Der Zaubernde ruft hiermit alle Vögel innerhalb eine Radius von \textbf{Stärke}×100 Metern herbei, die der Meinung sind am gewünschten Ort würde sich eine große Menge ihres Lieblingsfutters befinden. Der Zaubernde kann die Vögel nicht kontrollieren, aber ein geübter Beobachter ist in der Lage die einzelnen Vogelarten zu erkennen.

Magieniveau 4+: Der Zaubernde kann den Vögeln nahe legen, eine Aufgabe für ihn zu erfüllen. Gelingt ein Wurf auf Charme mit \textbf{Magieniveau} Erfolgen, so kann es sein dass sie ihn erhören.

\spellinfobox{3}{1}{0 Schritt}{---}{---}{1 Kampfhandlung}{Nein}


\subsection{Vertrauen}

\spellheader{Kontrolle}{Natur}{1 Kampfhandlung}
Das Tier fasst volles Vertrauen zu dem Zaubernden. Für \textbf{Stärke}×5 Minuten gilt das Tier als Vertrauter. Das Tier darf ein natürliches Lebewesen von der Größe eines Wolfes sein. Das Tier darf dem Zaubernden für diesen Zauber nicht feindlich gesonnen sein.

Magieniveau 4+: Das Tier darf die Größe eines Elefanten haben.

\spellinfobox{5}{1}{5 Schritt}{Stärke×5 Minuten}{---}{1 Kampfhandlung}{Nein}


\subsection{Tier finden}

\spellheader{Weissagung}{Natur}{Ritual}
Der Zaubernde begibt sich in den natürlichen Lebensraum eines Tieres (also z.B. an einen Bach oder auf einen Baum), und verharrt dort für 5 Minuten. Während dieser Zeit richtet er seine Gedanken an das Tier, welches er sucht.

Der Zaubernde spürt den genauen Aufenthaltsort des Tieres auf. Er weiss in dem Moment, wo sich das Tier befindet.

\spellinfobox{3}{1}{0 Schritt}{---}{---}{Ritual}{Nein}


\subsection{Salamanderfüße}

\spellheader{Transmutation}{Natur}{1 Kampfhandlung}
Der Zaubernde ist in der Lage, an glatten Wänden und sogar Decken oder Überhängen auf Händen und Füßen zu laufen.

Zudem kann er gefahrlos aus Höhen bis zu \textbf{Stärke}×3 Metern herunterspringen und landet unbeschadet auf den Füßen. Er braucht jedoch freie Hände, um sich damit abzufangen. Schuhe und Handschuhe kann er hierfür anbehalten.

Der Effekt hält für \textbf{Magieniveau} Minuten an.

\spellinfobox{5}{1}{0 Schritt}{Magieniveau Minuten}{---}{1 Kampfhandlung}{Nein}


\subsection{Ruf der Wildnis}

\spellheader{Beschwörung}{Natur}{1 Kampfhandlung}
Der Zaubernde ruft eine Gruppe einer gewünschten Tierart aus bis zu \textbf{Stärke} Kilometern Umkreis zusammen. Die Tiere verhalten sich für die Dauer des Zaubers wie Vertraute des Zaubernden.

Nach Ablauf des Zaubers wirft der Zaubernde auf Charme. Misslingt der Wurf, richten sich die Tiere gegen den Zaubernden oder seine Gefährten.

Der Zaubernde kann bis zu \textbf{Stärke} Tiere von der Größe eines Wolfes rufen.

\spellinfobox{3}{3}{0 Schritt}{Magieniveau Minuten}{---}{1 Kampfhandlung}{Nein}


\subsection{Reines Wasser}

\spellheader{Transmutation}{Natur}{1 Kampfhandlung}
Das Wasser, welches der Zaubernde berührt, wird, ausgehend von der Hand des Zaubernden so rein, dass man es trinken kann. Die endgültigen Kosten hängen dabei von der Reinheit des Wassers vor dem Wirken des Zaubers ab. Der Spielleiter entscheidet darüber. Der Zaubernde kann bis zu \textbf{Stärke} Eimer Wasser reinigen.

Magieniveau 4+: Das Wasser erhält eine heilende Kraft. Davon zu trinken heilt \textbf{Magieniveau} Wunden.

\spellinfobox{2}{1}{0 Schritt}{---}{---}{1 Kampfhandlung}{Nein}


\subsection{Rauch}

\spellheader{Beschwörung}{Feuer}{1 Kampfhandlung}
Der Zaubernde wirft ein brennendes Stück Stoff auf den Boden.

Von dem brennenden Stück Stoff geht, sobald es auf dem Boden aufgeprallt ist, ein dichter Rauch aus. Der Rauch wird vom Wind getrieben und ist nicht kontrollierbar. Es entsteht genug in den Augen brennender Rauch um eine Wolke von \textbf{Stärke}×10 Metern um den Zaubernden zu bilden.

\spellinfobox{5}{1}{0 Schritt}{Magieniveau Minuten}{Wolke}{1 Kampfhandlung}{Nein}


\subsection{Tier markieren}

\spellheader{Weissagung}{Natur}{1 Kampfhandlung}
Der Zaubernde markiert das Tier, so dass er über den ganzen Zeitraum des Zaubers die Richtung kennt, in der sich das Tier befindet. Der Zauber wirkt \textbf{Stärke} Tage. Der Zaubernde ist in der Lage, die Entfernung des Tieres zu bestimmen.

\spellinfobox{2}{1}{0 Schritt}{Stärke Tage}{---}{1 Kampfhandlung}{Nein}


\subsection{Lebendiges Versteck}

\spellheader{Transmutation}{Natur}{1 Kampfhandlung}
Der Zaubernde lässt einen Gegenstand, den er verstecken möchte, in eine Pflanze gleiten.

Der Gegenstand dringt in die Pflanze ein, ohne sie dabei zu beschädigen. Um den Gegenstand wieder aus der Pflanze zu befreien, muss der Zaubernde den Zauber erneut ausführen. Der Gegenstand kann \textbf{Stärke} Monde in der Pflanze bleiben, danach wird er von ihr wieder ausgeworfen.

\spellinfobox{5}{3}{0 Schritt}{Stärke Monate}{---}{1 Kampfhandlung}{Nein}


\subsection{Heilende Hände}

\spellheader{Heilung}{Natur}{1 Kampfhandlung}
Der Zaubernde heilt \textbf{Stärke}+\textbf{Magieniveau} Wunden beim Verwundeten. Dabei werden auch Knochenbrüche geheilt und Wunden verschlossen.

\spellinfobox{5}{1}{0 Schritt}{---}{---}{1 Kampfhandlung}{Nein}


\subsection{Falkenaugen}

\spellheader{Verzauberung}{Natur}{1 Kampfhandlung}
Der Zaubernde kann für eine Stunde alles bis in eine Entfernung von \textbf{Stärke} Kilometern klar und genau sehen, es sei denn, Nebel oder Rauch behindern seine Sicht. Angriffe mit Fernkampfwaffen erhalten \textbf{Magieniveau} Trefferwürfel zusätzlich.

\spellinfobox{5}{1}{0 Schritt}{1 Stunden}{---}{1 Kampfhandlung}{Nein}


\subsection{Naturspiel}

\spellheader{Verzauberung}{Natur}{1 Kampfhandlung}
Dem Zaubernden und den anderen Zusachauern zeigt sich die Natur in ihrer ganzen Anmut, Schönheit und Perfektion. Eine perfekte Landschaft offenbart sich ihnen: Die Bäume strahlen in ihren schönsten Farben und wippen rythmisch zum angenehmen Säuseln des Windes. Wasser wirft prächtige Wellenspiele, über welches Fisch symmetrisch hinwegspringen. Das Naturspiel ist in seiner Ausprägung jedesmal anders und natürlich auch abhnängig von der Gessinnung des Zaubernden.

Der Zauber dauert \textbf{Stärke}+\textbf{Magieniveau} Minuten an.

\spellinfobox{3}{1}{0 Schritt}{Stärke+Magieniveau Minuten}{---}{1 Kampfhandlung}{Nein}


\subsection{Tiergefährte}

\spellheader{Verzauberung}{Natur}{1 Kampfhandlung}
Der Zaubernde stellt eine geistige Verbindung zu einem Tier in Reichweite her. Das Tier braucht Erfolge anhand der \textbf{Stärke} auf seine Willenskraft oder ist für die nächsten \textbf{Magieniveau} Stunden sein treuer Gefährte.

Wenn dem Zaubernden das Tier feindlich gesonnen ist ist der Magiewurf eine schwere Probe (+1).

\spellinfobox{5}{1}{50 Schritt}{8 Stunden}{---}{1 Kampfhandlung}{Nein}


\subsection{Bioluminiszenz}

\spellheader{Beschwörung}{Erde}{1 Kampfhandlung}
Der Zaubernde erschafft \textbf{Stärke} lebende Vorkommen von Bioluminiszenz im Umkreis von 10 Schritt. Diese Wesen bestehen aus lebendiger pflanzlicher Materie. Sie beleuchten die Umgebung auf magische Weise, und können sich sehr langsam bewegen (Schnelligkeit 1).

Die Wesen folgen dem Zaubernden und sorgen für die Dauer ihrer Existenz für ein angenehmes Licht um den Zaubernden herum. Sie haben keine Kampffähigkeiten, es sind Pflanzen, die über den Boden kriechen.

Sie bleiben \textbf{Magieniveau} W6 Stunden bestehen, dann zerfallen sie.

\spellinfobox{5}{1}{0 Schritt}{Magieniveau W6 Runden}{---}{1 Kampfhandlung}{Nein}


\section{Schwarze Magie}

\subsection{Alp rufen}

\spellheader{Beschwörung}{Geist}{1 Kampfhandlung}
Der Zaubernde ruft einen Alp, ein Geisterwesen, welches seine Opfer im Schlaf heimsucht. Der Alp hat keine Gestalt, und kann beliebig und ohne Zeitverlust den Ort wechseln. Der Geist unterliegt nicht der Beherrschung des Zaubernden. Soll der Alp einen Dienst erweisen, so ist ein gelungener Charme-Wurf erforderlich. Misslingt dieser, wird der Alp so schnell wie möglich verschwinden.

Der Mindestwurf des Charme-Wurfs ist 7-\textbf{Magieniveau}. Der Zaubernde kann dem Alp \textbf{Stärke} Befehle geben.

\spellinfobox{5}{2}{0 Schritt}{1 Nächte}{---}{1 Kampfhandlung}{Nein}


\subsection{Beherrschung}

\spellheader{Kontrolle}{Geist}{1 Kampfhandlung}
Der Zaubernde kann das Opfer dazu bringen, einen Befehl des Zaubernden auszuführen.

Der vom Zaubernden übertragene Befehl manifestiert sich als "fixe Idee" im Geist des Beherrschten. Diese drängt sich während der Wirkungsdauer des Zaubers immer wieder in die Gedanken des Beherrschten, solange dieser ihr nicht nachgibt, und kann so bisweilen sogar seine Konzentration auf andere Dinge unangenehm stören.

Der Zaubernde kann den Beherrschten dazu bringen, Dinge zu tun, die für ihn nicht unmittelbar lebensbedrohlich sind oder gegen seine innersten Prinzipien verstoßen. So wird man etwa einen Magier nicht dazu bringen können, wertvolle Bücher zu verbrennen, wohl aber einen bezahlten Handlanger, seinen Meister zu verraten oder einen Söldner, im Kampf die Seiten zu wechseln.

Der Zauber dauert \textbf{Stärke} Minuten an. Einmal pro Minute wirft das Opfer auf seine Willenskraft. Erziehlt es weniger Erfolge als  \textbf{Stärke des Zaubers}+\textbf{Magieniveau}, so ist es von der fixen Idee überzeugt und führt sie aus.

\spellinfobox{7}{2}{10 Schritt}{Stärke×5 Minuten}{---}{1 Kampfhandlung}{Ja}


\subsection{Blutdornen}

\spellheader{Beschwörung}{Blut}{1 Kampfhandlung}
Der Zauberer lässt auf einer vorher festgelegten, \textbf{Stärke}×2 Meter großen, runden Fläche dunkle Dornen aus dem Boden sprießen. Jeder, der sich am Ende seiner eigenen Kampfrunde darin befindet, erhält  \textbf{Stärke} Treffer.

Die Blutdornen bleiben für \textbf{Magieniveau}+1 Kampfrunden bestehen und verfallen danach zu einem schwarzen, stinkenden Brei.

\spellinfobox{7}{2}{20 Schritt}{Magieniveau+1 Runden}{Kreis}{1 Kampfhandlung}{Nein}


\subsection{Blutfluss}

\spellheader{Transmutation}{Blut}{1 Kampfhandlung}
Das Opfer bäumt sich auf, während ein Schwall Blut blitzschnell durch alle Poren aus seinem Körper sickert. Pro Kampfrunde nimmt das Opfer \textbf{Magieniveau} Wunden hin, der Zauber dauert \textbf{Stärke} Kampfrunden an.

Am Ende der Kampfrunde wirft das Opfer auf seine Kraft. Erreicht es so viele Erfolge wie der Zauber \textbf{Stärke} hat, so endet der Zauber sofort.

\spellinfobox{7}{2}{0 Schritt}{Stärke Runden}{---}{1 Kampfhandlung}{Ja}


\subsection{Dunkle Pfade}

\spellheader{Kontrolle}{Geist}{1 Kampfhandlung}
Der Magier kann sich im Umkreis von (\textbf{Stärke}+\textbf{Magieniveau})×3 Metern nach seinem Willen translokieren. Während der Translokation scheint ein dunkler Schatten über den Boden zu gleiten.

\spellinfobox{7}{1}{0 Schritt}{---}{---}{1 Kampfhandlung}{Nein}


\subsection{Eins mit den Schatten}

\spellheader{Verzauberung}{Licht}{1 Kampfhandlung}
Solange der Verzauberte im Schatten bleibt, ist er nur schwer auszumachen. Er ist dann selbst nur als tieferer Schatten innerhalb des Schattens zu erkennen.

Wenn er stillsteht, werden ihn nur Personen entdecken, die gezielt nach ihm suchen und direkt auf ihn schauen. Hierfür müssen bei einem Wahrnehmungs-Wurf \textbf{Stärke} Erfolge erreicht werden.

Bewegt sich der Verzauberte langsam, ist die Anzahl der zu erreichenden Erfolge halbiert, während schnelle Bewegungen ohne Schwierigkeiten zu sehen sind.

Der Zauber währt \textbf{Stärke} Minuten.

\spellinfobox{5}{2}{0 Schritt}{Stärke Minuten}{---}{1 Kampfhandlung}{Nein}


\subsection{Nachtwesen}

\spellheader{Verzauberung}{Licht}{1 Kampfhandlung}
Sobald die Sonne untergegangen ist gibt der Verzauberte beim Gehen kaum einen Laut von sich (Die Fertigkeit Heimlichkeit ist um \textbf{Stärke}+\textbf{Magieniveau} des Zaubers erhöht) und kann im Dunkeln wie bei Tage sehen.

Der Zauber dauert eine Nacht an.

\spellinfobox{5}{2}{0 Schritt}{1 Nächte}{---}{1 Kampfhandlung}{Nein}


\subsection{Stimme des Hasses}

\spellheader{Kontrolle}{Geist}{1 Kampfhandlung}
Der Zaubernde flüstert dem ersten Opfer irgendetwas ins Ohr.

Das Opfer wird von Hass durchzogen und empfindet das Verlangen, den Zauber weiterzutragen, indem es bis zu \textbf{Stärke} weitere Opfer anspricht. So baut sich den Hass immer weiter auf.

Der Effekt dauert \textbf{Magieniveau} Tage an.

\spellinfobox{11}{3}{0 Schritt}{Magieniveau Tage}{---}{1 Kampfhandlung}{Nein}


\subsection{Vision der Angst}

\spellheader{Verzauberung}{Geist}{1 Kampfhandlung}
Das Opfer empfindet grausame Bilder verzerrter Welten, die es als Vision empfindet. Mehr und mehr kommt es ihm jedoch als Realität vor. Misslingt ein Willenskraft Wurf gegen den Mindestwurf \textbf{Magnieniveau}+2, droht das Opfer kurzzeitig wahnsinnig zu werden.

Der Zauber dauert \textbf{Stärke} Minuten an.

\spellinfobox{7}{1}{0 Schritt}{Stärke Minuten}{---}{1 Kampfhandlung}{Nein}


\subsection{Zungen der Seth'Nra}

\spellheader{Beschwörung}{Dämonisch}{1 Kampfhandlung}
In einem Kreis von \textbf{Stärke} Meter Durchmesser erwachsen \textbf{Stärke} tentakelähnliche schwarze Zungen aus dem Boden, die Lebewesen jeglicher Art festhalten. Durch die Berührung der Tentakel sinkt der Arkanawert der Festgehaltenen um \textbf{Magieniveau} Punkte pro Kampfrunde. Physischer Schaden entsteht nicht. Die Tentakel haben 3 mögliche Wunden. Werden sie nicht vernichtet, so bleiben sie auf unbestimmte Zeit bestehen.

\spellinfobox{5}{2}{10 Schritt}{---}{Kreis}{1 Kampfhandlung}{Nein}


\subsection{Schwarzer Tod}

\spellheader{Schaden}{Blut}{1 Kampfhandlung}
Der Zaubernde rezitiert die folgenden Worte mit klarer, fester Stimme:

"Aus der Tiefe der Schatten, aus dem Schoß des Todes, rufe ich die schwarze Pest. Faul sei das Fleisch, welk sei der Geist, bis die Stille des Grabes dich heimführt!"

Schwarze, geisterhafte Schwaden steigen aus deiner Hand empor und rasen auf das Ziel zu. Das Opfer verspürt sofort eine brennende Kälte in seinen Adern. Innerhalb von Sekunden breiten sich dunkle Flecken auf seiner Haut aus, begleitet von starkem Schwindel und Schwäche.

Falls der Zauber nicht rechtzeitig gebannt oder geheilt wird, verfällt das Ziel innerhalb von Minuten einem tödlichen Fieber.

Der Zauber fügt direkten Schaden zu und ignoriert jeglichen Rüstungswert des Ziels.

Wirft der Zaubernde einen Megakrit zerfällt das Ziel sofort bei Kontakt des Zaubers zu Staub.

Höhere Wesen erleiden durch diesen Zauber nicht den sofortigen Tod. Diese erleiden im Verlauf des Kampfes pro Runde Magieniveau × Stärke Schaden.

Gegenmittel: Eine Priestersegnung oder ein starker Heilzauber kann den Fluch aufheben, wenn er innerhalb einer Stunde gewirkt wird. Das Blut eines Engels kann den Effekt vollständig neutralisieren.

\spellinfobox{20}{10}{20 Schritt}{---}{Strahl}{1 Kampfhandlung}{Nein}


\section{Verfluchungen}

\subsection{Altern}

\spellheader{Transmutation}{Blut}{1 Kampfhandlung}
Der Zaubernde muss den Fluch sprechen und das Opfer dabei berühren.

Kurz nach der Verfluchung beginnnen beim Opfer die ersten Auswirkungen einzusetzen. Der Verfluchte fühlt sich geschwächt, dies äussert sich neben dem unangenehmen Gefühl in der Reduzierung der Werte für Kraft, Willenskraft, Reaktion, Schnelligkeit und Wahrnehmung um 1.

So lange der Fluch aktiv ist, verringern sich diese Werte jeden Tag um einen weiteren Punkt. Erreicht eines der Attribute den Wert 0, so wird das Opfer bettlägerig und kann nicht mehr von alleine aufstehen.

Der Fluch kann durch Gegenmagie, göttliches Wirken oder einen Segen gebrochen werden. Ebenso endet der Fluch, wenn der Zaubernde den Zauber fallen lässt oder ohnmächtig wird.

\spellinfobox{5}{2}{0 Schritt}{Stärke+Magieniveau Jahre}{---}{1 Kampfhandlung}{Ja}


\subsection{Angst}

\spellheader{Kontrolle}{Geist}{1 Kampfhandlung}
Der Zaubernde muss den Fluch sprechen und das Opfer dabei berühren.

Ein Opfer verfällt in eine regelrechte Panik und denkt für die nächsten \textbf{Magieniveau} W6 Minuten nurnoch an die Flucht.  Das Opfer dem Fluch durch einen gelungenen Tapferkeitswurf entgehen. Der Mindestwurf dieses Wurfs ist um die Zauberstärke des Zaubers erhöht.

Der Fluch kann durch Gegenmagie, göttliches Wirken oder einen Segen gebrochen werden. Ebenso endet der Fluch, wenn der Zaubernde den Zauber fallen lässt oder ohnmächtig wird.

\spellinfobox{5}{2}{0 Schritt}{Magieniveau W6 Minuten}{---}{1 Kampfhandlung}{Nein}


\subsection{Austrocknen}

\spellheader{Transmutation}{Blut}{1 Kampfhandlung}
Der Zaubernde muss den Fluch sprechen und das Opfer dabei berühren.

Das Opfer verliert jeden Tag \textbf{Magieniveau} + 5\% seiner Körperflüssigkeit. Durch Aufnahme von Flüssigkeit kann er den Effekt auf 3\% verlangsamen, jedoch nicht ganz verhindern. Nachdem der Verfluchte 20\% seiner Körperflüssigkeit verloren hat, sieht man bereits erste Anzeichen des Mangels. Die Haut wird trockener, es bilden sich erste wunde Stellen. Bei 40\% ist der Verfluchte bereits so weit geschwächt, dass alle Attributwerte halbiert werden. Ab einem Flüssigkeitsverlust von 70\% (jetzt scheint die Haut nunmehr einem trockenen Stück Leder zu gleichen) ist es dem Verfluchten kaum mehr möglich, sich aufzurichten. Alles um ihn herum verschwimmt, und er hat keine Möglichkeit, sich selbstständig zu versorgen. Erst wenn alle Flüssigkeit aus dem Körper entwichen ist, stirbt der Verfluchte. Bis zu diesem Zeitpunkt hält ihn der Fuch am Leben.

Der Fluch kann durch Gegenmagie, göttliches Wirken oder einen Segen gebrochen werden. Ebenso endet der Fluch, wenn der Zaubernde den Zauber fallen lässt oder ohnmächtig wird.

\spellinfobox{5}{3}{0 Schritt}{Stärke Jahre}{---}{1 Kampfhandlung}{Nein}


\subsection{Fluch des Geistes}

\spellheader{Schaden}{Blut}{Ritual}
Der Zaubernde muss den Fluch sprechen und dabei den Namen des Opfers auf ein Blatt Papier oder Pergament schreiben.

Der Zaubernde begibt sich in einen Trance ähnlichen Zustand, in dem er das vorzugsweise betäubte Opfer ausbluten lässt. Der Fluch wird auf den übertragen, dessen Namen auf dem Papier geschrieben steht, und auf den der Zaubernde seinen Hass fokusiert hat. Danach fällt der Zaubernde in eine Ohnmacht die in einen unruhigen Schlaf übergleitet.

Das Opfer blutet fortan, so lange der Fluch anhält, aus allen Poren. Blut tritt aus seinen Augen und seiner Nase, und das Opfer nimmt \textbf{Magieniveau} W6 Wunden pro Tag hin.

Der Fluch kann durch Gegenmagie, göttliches Wirken oder einen Segen gebrochen werden. Ebenso endet der Fluch, wenn der Zaubernde den Zauber fallen lässt oder ohnmächtig wird.

\spellinfobox{5}{3}{40 Schritt}{Stärke Monate}{---}{Ritual}{Ja}


\subsection{Lähmung des Selbst}

\spellheader{Kontrolle}{Geist}{1 Kampfhandlung}
Der Zaubernde muss den Fluch sprechen und das Opfer dabei berühren.

Die Zunge und Glieder des Opfers fühlen sich schwer an und möchten nicht mehr richtig ihren Dienst tun.  Alle körperlichen Attribute ausser Resistenz reduzieren sich mit allen Konsequenzen um 1. Der Fluch wirkt \textbf{Stärke}+\textbf{Magieniveau} Stunden.

Der Fluch kann durch Gegenmagie, göttliches Wirken oder einen Segen gebrochen werden. Ebenso endet der Fluch, wenn der Zaubernde den Zauber fallen lässt oder ohnmächtig wird.

\spellinfobox{5}{2}{0 Schritt}{Stärke+Magieniveau Stunden}{---}{1 Kampfhandlung}{Nein}


\subsection{Madaeus Grippe}

\spellheader{Transmutation}{Blut}{1 Kampfhandlung}
Der Zaubernde muss den Fluch sprechen und das Opfer dabei berühren.

Das Opfer erkrankt am nächsten Tag an der Madeus Grippe. Hals und Rachenbeschwerden, neben leichtem Fieber und Hustenreizen kennzeichnen das Krankheitsbild. Das Opfer bleibt solange krank bis er entweder auf magischem Wege geheilt wird, oder mindestens zwei weitere Personen auf natürlichem Wege sich bei ihm angesteckt haben. Die Grippe verläuft nie tödlich, wird jedoch als lästiges, unangenehmes Ärgernis empfunden. Die "Ansteckungsrate" beträgt \textbf{Magieniveau}×10 \% pro Tag Aufenthalt in Begleitung eines Verfluchten.

Der Fluch kann durch Gegenmagie, göttliches Wirken oder einen Segen gebrochen werden. Ebenso endet der Fluch, wenn der Zaubernde den Zauber fallen lässt oder ohnmächtig wird.

\spellinfobox{5}{1}{0 Schritt}{---}{---}{1 Kampfhandlung}{Nein}


\subsection{Pech}

\spellheader{Kontrolle}{Geist}{1 Kampfhandlung}
Der Zaubernde muss den Fluch sprechen und das Opfer dabei berühren.

Das Opfer des Fluchs erhält für die Dauer des Fluchs einen Wert "Pech" in Höhe der Stärke des Zaubers.

Das Opfer muss, nachdem es eine Probe für seine Handlungen vorgenommen hat, auf den "Pech" Wert würfeln. Zeigt der "Pech" Wurf einen Erfolg, so misslingt die Handlung.

Der Fluch kann durch Gegenmagie, göttliches Wirken oder einen Segen gebrochen werden. Ebenso endet der Fluch, wenn der Zaubernde den Zauber fallen lässt oder ohnmächtig wird.

\spellinfobox{5}{1}{0 Schritt}{Stärke+Magieniveau Stunden}{---}{1 Kampfhandlung}{Nein}


\subsection{Selbst Schuld}

\spellheader{Kontrolle}{Geist}{1 Kampfhandlung}
Der Zaubernde muss den Fluch sprechen und das Opfer dabei berühren.

Das Opfer spürt nichts von der Verfluchung bis es den ersten Schlag/Schuss gegen irgendein Ziel ausführt.

Der Schaden der bei einem Treffer vom Verfluchten an einem Ziel verursacht wird kommt 1 zu 1 zu ihm zurück, hierbei erleidet er die gleichen Treffer wie der Angegriffene. Der Fluch dauert \textbf{Magieniveau} Angriffe vom Verfluchten an.

Der Fluch kann durch Gegenmagie, göttliches Wirken oder einen Segen gebrochen werden. Ebenso endet der Fluch, wenn der Zaubernde den Zauber fallen lässt oder ohnmächtig wird.

\spellinfobox{5}{1}{0 Schritt}{---}{---}{1 Kampfhandlung}{Nein}


\subsection{Shuras Wahn}

\spellheader{Kontrolle}{Geist}{1 Kampfhandlung}
Der Zaubernde muss den Fluch sprechen und das Opfer dabei berühren.

Das Opfer wird unmittelbar von panischer Angst ergriffen. Der einzige Gedanke, der für die nächsten \textbf{Magieniveau} W6 Sekunden gefasst werden kann ist "WEG HIER". Allerdings kann die Angst durch eine Willenskraftprobe abgeschüttelt werden. Der Mindestwurf dieser Probe ist um die Stärke des Zaubers erhöht.

Der Fluch kann durch Gegenmagie, göttliches Wirken oder einen Segen gebrochen werden. Ebenso endet der Fluch, wenn der Zaubernde den Zauber fallen lässt oder ohnmächtig wird.

\spellinfobox{5}{1}{0 Schritt}{Magieniveau W6 Sekunden}{---}{1 Kampfhandlung}{Nein}


\section{Weiße Magie}

\subsection{Wehrlose Gestalt}

\spellheader{Illusion}{Arkana}{1 Kampfhandlung}
Der Zaubernde wirkt für 5×\textbf{Stärke} Minuten absolout harmlos. Je nach Erscheinung erscheint er wie ein gebrechlicher alter, kranker Mann, eine wehrlose Frau oder Ähnlichem.

Der Zaubernde erhält in der Zeit einen Bonus von \textbf{Magieniveau} Punkten auf seine \textit{Heimlichkeit}.

\spellinfobox{5}{1}{0 Schritt}{5×Stärke Minuten}{---}{1 Kampfhandlung}{Nein}


\subsection{Überlicht}

\spellheader{Kontrolle}{Licht}{1 Kampfhandlung}
Dem Zaubernden ist es möglich ca. \textbf{Stärke}×2 Sekunden stattgefundenen Wirkungen in einer Reichweite von \textbf{Magieniveau}×5 Schritt zuvor zu kommen. Es ist ihm so möglich z.B. ein Glas noch zu fangen obwohl es bereits am Boden zerschellt.

\spellinfobox{5}{2}{15 Schritt}{---}{---}{1 Kampfhandlung}{Nein}


\subsection{Magier erkennen}

\spellheader{Weissagung}{Arkana}{1 Kampfhandlung}
Der Zaubernde kann die magische Begabung und die Ausrichtung einer Person in seinem Sichtfeld erkennen. Der Beobachtete wirft einen Willenskraft Wurf. Erreicht er Erfolge entsprechend der \textbf{Stärke des Zaubers} + \textbf{Magieniveau}, so bleibt seine magische Begabung verborgen.

\spellinfobox{3}{1}{200 Schritt}{---}{---}{1 Kampfhandlung}{Nein}


\subsection{Schutzwand}

\spellheader{Beschwörung}{Erde}{1 Kampfhandlung}
Der Zaubernde erzeugt um sich herum eine Schutzwand die dem Zaubernden Schutz bietet.

Die Wand hält \textbf{Stärke}×2 Wunden aus. Die Schutzwand kann nur durch magische Waffen oder Zauber beschädigt werden, auf diese Weise kann sie vorzeitig auf 0 abgetragen werden.

Die Wand bleibt \textbf{Magieniveau}+1 Kampfrunden bestehen.

\spellinfobox{5}{1}{0 Schritt}{Stärke Runden}{---}{1 Kampfhandlung}{Nein}


\subsection{Niedere Untote bannen}

\spellheader{Widerrufung}{Licht}{1 Kampfhandlung}
Aus der Hand des Zaubernden entfährt ein grellend weisser Lichtblitz, der auf bis zu \textbf{Stärke} Untote einschlägt. Diese zerbersten sofort und hinterlassen nichts als einen rauchenden Haufen Gebeine. Der Zauber wirkt nur bei niederen Untoten wie Zombies, Skeletten, oder niederen Vampiren. Höheren Untoten wie z.B. höhere Vampire oder Werwesen fügt er erheblichen Schaden zu (\textbf{Magieniveau}×3 Wunden).

\spellinfobox{9}{3}{10 Schritt}{---}{---}{1 Kampfhandlung}{Nein}


\subsection{Magie bannen}

\spellheader{Widerrufung}{Arkana}{1 Kampfhandlung}
Der Zaubernde erzeugt einen \textbf{Stärke}×2 Meter großen Magiebannkreis. Keine Magie oder magische Handlung kann in diesem Umkreis gewirkt werden. Bestehende Zauber verfallen sofort, mit ausnahme des Zaubers "Magie bannen".

Der Kreis bleibt \textbf{Magieniveau}+1 Minuten bestehen.

\spellinfobox{3}{1}{0 Schritt}{5 Minuten}{Kreis}{1 Kampfhandlung}{Nein}


\subsection{Lüge erkennen}

\spellheader{Weissagung}{Geist}{1 Kampfhandlung}
Der Zaubernde erkennt ob sein Gegenüber lügt oder nicht. Es können bis zu \textbf{Magieniveau} Aussagen des Beobachteten geprüft werden.

Das Ziel des Zaubers darf auf seine Willenskraft werfen. Gelingt der Wurf mit \textbf{Stärke} Erfolgen, so bleibt verborgen ob es lügt oder nicht.

\spellinfobox{5}{1}{3 Schritt}{---}{---}{1 Kampfhandlung}{Nein}


\subsection{Lichtkreis}

\spellheader{Beschwörung}{Licht}{1 Kampfhandlung}
Der Zaubernde erzeugt einen hellen Lichtkreis von \textbf{Stärke}×2 Metern Umkreis um sich herum. Der Kreis bleibt 15 Minuten bestehen.

\spellinfobox{3}{1}{0 Schritt}{15 Minuten}{Kreis}{1 Kampfhandlung}{Nein}


\subsection{Lichtgeschwind}

\spellheader{Transmutation}{Licht}{1 Kampfhandlung}
Der Zaubernde beginnt mit einer Geschwindigkeit gleich der des Lichtes zu laufen. Für Umstehende scheint er sich in Luft aufzulösen, in Wahrheit aber sprintet er mit unfassbarer Geschwindigkeit, nimmt seine Umgebung aber wahr als würde er lediglich schnell laufen. Es ist dem Zaubernden möglich über jegliches, begehbares Terrain zu laufen, so kann er beispielsweise einen kompletten Kontinent in der Zeit eines Augenschlages durchqueren. Es ist ihm nicht möglich während des Laufens eine andere Handlung zu vollziehen.

Der Zauber wirkt eine Sekunde lang.

\spellinfobox{11}{3}{0 Schritt}{---}{---}{1 Kampfhandlung}{Nein}


\subsection{Lichtattacke}

\spellheader{Schaden}{Licht}{1 Kampfhandlung}
Das Opfer wird von gleissend hellem Licht geblendet und ist für \textbf{Stärke} Kampfrunden vollkommen orientierungslos und zu keiner Handlung in der Lage. Das Opfer nimmt \textbf{Magieniveau} Treffer hin.

\spellinfobox{5}{1}{0 Schritt}{Stärke Runden}{---}{1 Kampfhandlung}{Nein}


\subsection{Kerze}

\spellheader{Beschwörung}{Feuer}{1 Kampfhandlung}
Der Zaubernde entzündet bis zu \textbf{Stärke}×3 Kerzen.

\spellinfobox{3}{1}{50 Schritt}{---}{---}{1 Kampfhandlung}{Nein}


\subsection{Höhere Untote bannen}

\spellheader{Widerrufung}{Licht}{1 Kampfhandlung}
Aus den Händen des Zaubernden entfährt eine grellend weisse Lichtwand, die auf einen Untoten zugleitet. Dieser lodert in Flammen auf und erleidet dabei unvorstellbare Qualen (man sagt er durchlebe alle Qualen seiner Opfer auf einmal). Ausser einem Haufen Asche bleibt von dem Höheren Untoten nichts mehr übrig. Die Lichtwand hat etwa eine Breite von \textbf{Stärke}×3 Metern, es ist auch möglich mehrere niedere Untote zu vernichten, die von der Wand getroffen werden, jedoch nur einen höheren Untoten. Weitere höhere Untote, die sich neben dem Opfer befinden, erleiden \textbf{Magieniveau}×2 Wunden.

Die Wand bewegt sich mit einer Geschwindigkeit von \textbf{Magieniveau}+1 Schritt pro Kampfrunde.

\spellinfobox{13}{5}{0 Schritt}{---}{Wand}{1 Kampfhandlung}{Nein}


\subsection{Guter Freund}

\spellheader{Kontrolle}{Geist}{1 Kampfhandlung}
Der Zaubernde lässt das Opfer glauben, dass er ein guter Freund von ihm ist, ja sogar einer seiner Besten. Er erzählt im bereitwillig alles, was er auch seinem besten Freund erzählen würde. Nach Beenden des Zaubers kann das Opfer sich nicht erklären warum er dies getan hat.

Das Opfer des Zaubers wirft auf seine Willenskraft. Erreicht er Erfolge in Höhe der \textbf{Stärke des Zaubers}, so ist der Zauber fehlgeschlagen, und das Opfer hat Kenntnis von dem Versuch der Verzauberung.

Die Freundschaft hält \textbf{Maginiveau} Minuten an.

\spellinfobox{7}{2}{0 Schritt}{Magieniveau W6 Minuten}{---}{1 Kampfhandlung}{Nein}


\subsection{Geister bannen}

\spellheader{Widerrufung}{Arkana}{1 Kampfhandlung}
Der Zaubernde bannt  bis zu \textbf{Magieniveau} Geisterwesen, die sich in einem definiertem Gebiet aufhalten (Haus/Tempel/Waldstück). Er muss dabei die Geisterwesen zumindestens gedanklich fixieren.

Die Geister werfen Würfel entsprechend ihrer verbleibenden Wunden. Erreichen sie so viele Erfolge wie die \textbf{Stärke des Zaubers} ist, so bleiben sie von dem Bann unberührt.

\spellinfobox{5}{3}{0 Schritt}{---}{---}{1 Kampfhandlung}{Nein}


\subsection{Furchteinflößende Gestalt}

\spellheader{Illusion}{Geist}{1 Kampfhandlung}
Der Zaubernde erscheint vor den Umstehenden als furchteinflössender Magier. Blitze zucken um ihn herum und Wind lässt seine Kleider aufwallen. Jeder der vorhatte sich dem Begabten zu nähern muss einen \textit{Tapferkeits}-Wurf mit \textbf{Stärke} Erfolgen bestehen.

\spellinfobox{3}{1}{0 Schritt}{Magieniveau+1 W6 Minuten}{---}{1 Kampfhandlung}{Ja}


\subsection{Verschwimmen}

\spellheader{Transmutation}{Licht}{1 Kampfhandlung}
Der Zauberne bricht das Licht um sich und lässt sein Form verschwimmen. Für \textit{Zauberstärke} Runden sind Angriffe gegen ihn erschwert (Mindestwurf + \textbf{Magieniveau}).

\spellinfobox{5}{1}{1 Schritt}{Stärke Runden}{---}{1 Kampfhandlung}{Nein}


\subsection{Flüche Bannen}

\spellheader{Widerrufung}{Geist}{1 Kampfhandlung}
Der Zaubernde bannt einen Fluch. Der Mindestwurf der Probe ist um die Stärke des Fluch Zaubers erschwert, und um das \textbf{Magieniveau} verringert.

\spellinfobox{4}{1}{10 Schritt}{---}{---}{1 Kampfhandlung}{Nein}


\subsection{Heilung}

\spellheader{Heilung}{Blut}{1 Kampfhandlung}
Der Zaubernde heilt das Ziel für \textbf{Stärke}×\textbf{Magieniveau} Wunden.

\spellinfobox{10}{2}{0 Schritt}{---}{---}{1 Kampfhandlung}{Nein}


\subsection{Neronstatue}

\spellheader{Illusion}{Licht}{1 Kampfhandlung}
Der Begabte läßt eine transparente Statue erscheinen, die nach dem Vorbild - einer Statue des jungen Neron im Zentrum der Neronitenniederlassung - geformt ist. Dadurch dass nie ein Neroniete den lebendigen Neron gesehen hat, wirkt sie auch dementsprechend statisch und leblos. Je nach Willen des Zaubernden hat die Statue eine Größe von 10cm bis lebensgroß.

\spellinfobox{3}{1}{0 Schritt}{Magieniveau Stunden}{---}{1 Kampfhandlung}{Nein}


\subsection{Schwarzmagus erkennen}

\spellheader{Weissagung}{Geist}{1 Kampfhandlung}
Der Magus kann bis zu \textbf{Magieniveau} W6 Tage die Gesinnung, wie auch Ausrichtung eines jeden Magiers, den er sieht, erkennen. Versperrt sich ein Magier dieser Untersuchung durch seine Magiekunde, so durchleuchtet der Weissmagier immer noch seinen Gesinnungswert. Der Spruch ist vor allem gegen sein schwarzmagisches Gegenstück Weissmagus konzipiert, er neutralisiert diesen Zauber völlig.

\spellinfobox{3}{1}{0 Schritt}{Magieniveau W6 Tage}{---}{1 Kampfhandlung}{Nein}


\subsection{Verwandlung beenden}

\spellheader{Transmutation}{Blut}{1 Kampfhandlung}
Der Zaubernde bannt eine vor ihm stattfindende Verwandlung. Der Verwandelnde darf mit seinem Magiekunde Wert gegen den Bann des Zaubernden würfeln, wobei der Mindestwurf des Zaubers um die \textbf{Stärke} des "Verwandlung beenden" Zaubers erhöht ist. Zeigt der Wurf einen Erfolg, so bleibt die Verwandlung bestehen.

\spellinfobox{5}{1}{0 Schritt}{---}{---}{1 Kampfhandlung}{Nein}


\section{Zauberei}

\subsection{Bogus}

\spellheader{Illusion}{Arkana}{1 Kampfhandlung}
Der Zaubernde erzeugt eine Illusion, welche einen beliebigen Gegenstand ersetzt. Die Illusion muss etwa der Form des Gegenstands entsprechen. Der Gegenstand darf wie die Illusion eine Größe von Metern entsprechend der \textbf{Stärke des Zaubers} nicht übersteigen. Die Wirkungsdauer beträgt \textbf{Magieniveau}+1 W6 Minuten.

\spellinfobox{1}{1}{0 Schritt}{4W6 Minuten}{---}{1 Kampfhandlung}{Nein}


\subsection{Demaskieren}

\spellheader{Widerrufung}{Arkana}{1 Kampfhandlung}
Innerhalb eines Radius von \textbf{Stärke}+\textbf{Magieniveau} Metern hebt der Zauber alle Illusionen augenblicklich auf.

\spellinfobox{3}{1}{0 Schritt}{---}{Sphäre}{1 Kampfhandlung}{Nein}


\subsection{Glitzern}

\spellheader{Illusion}{Licht}{1 Kampfhandlung}
Im Sichtfeld des Zauberers entsteht an beliebiger Stelle ein Glitzern auf einer Fläche von 10×\textbf{Stärke} Zentimetern im Quadrat. Das Glitzern kann jede beliebige Form und Farbe annehmen. Das Glitzern bleibt für \textbf{Magieniveau} W6 Minuten bestehen.

\spellinfobox{3}{1}{200 Schritt}{Stärke W6 Minuten}{---}{1 Kampfhandlung}{Nein}


\subsection{Illusion}

\spellheader{Illusion}{Arkana}{1 Kampfhandlung}
Die große Illusion! Im Radius von \textbf{Stärke}×20 Schritt um den Zauberer beginnt sich alles zu verformen. Vertraute Gegenstände werden zu fremden Artefakten, Wände biegen sich zu unmöglichen Winkeln und Lebewesen wandeln sich zu andersartigen Kreaturen. Die Zone der Illusion bleibt an der Stelle, an der sie erzeugt wurde. Der Zauber währt (\textbf{Magieniveau}+1)W6 Minuten.

\spellinfobox{5}{2}{0 Schritt}{Stärke×2 W6 Minuten}{Sphäre}{1 Kampfhandlung}{Ja}


\subsection{Ogeratem}

\spellheader{Illusion}{Natur}{1 Kampfhandlung}
Der Zauberer erzeugt eine Wolke eines beliebigen Geruchs, die sich auf einer Fläche von \textbf{Stärke}×10 Schritt verbreitet. Der Geruch bleibt auch nach Beendigung des Zaubers bestehen, der Zauber ist jedoch \textbf{Magieniveau} W6 Minuten aktiv. Die Wolke lässt sich nur innerhalb der Zauberdauer steuern und wird danach vom Wind getrieben.

Maginiveau 5+: Der Geruch ist so intensiv, dass alle, die ihn riechen maßgeblich davon beeinflusst werden. Bei einem üblen Geruch wird ihnen übel, bei einem lieblichen Geruch wirkt es verzaubernd.

\spellinfobox{3}{1}{2 Schritt}{Stärke W6 Minuten}{Wolke}{1 Kampfhandlung}{Nein}


\subsection{Perfekte Form}

\spellheader{Illusion}{Arkana}{1 Kampfhandlung}
Der Zaubernde erscheint in einer perfekten Form. Alle guten Merkmale sind hervorgehoben. Der Zaubernde erhält für \textbf{Magieniveau} Minuten die \textbf{Stärke des Zaubers} als Bonus auf Attraktivität.

\spellinfobox{3}{1}{0 Schritt}{Magieniveau Minuten}{---}{1 Kampfhandlung}{Nein}


\subsection{Scheinbild}

\spellheader{Illusion}{Licht}{1 Kampfhandlung}
Der Zaubernde beschwört ein Scheinbild einer Kreatur, welches täuschen echt wirkt. Das Scheinbild bleibt für \textbf{Magieniveau} Runden bestehen. Ein erfolgreicher Wurf auf Aufassungsgabe gegen die \textbf{Zauberstärke} erlaubt es, dieses als Ilussion zu erkennen.

\spellinfobox{5}{2}{10 Schritt}{Magieniveau Runden}{---}{1 Kampfhandlung}{Nein}


\subsection{Doppelgänger}

\spellheader{Illusion}{Licht}{1 Kampfhandlung}
Der Zaubernde beschwört ein exaktes Abbild seiner selbst, das sich \textbf{Magieniveau} Kampfrunden lang in seiner Nähe bewegt und Angreifer verwirrt. Für die Wirkungsdauer des Zaubers wird der \textit{Ausweichwert} des Zaubernden um \textbf{Stärke} erhöht.

\spellinfobox{5}{2}{10 Schritt}{Magieniveau Runden}{---}{1 Kampfhandlung}{Nein}


\subsection{Wasser zu Wein}

\spellheader{Transmutation}{Wasser}{1 Kampfhandlung}
Der Zaubernde verwandelt Wasser, welches er in einem Gefäß vor sich hält, in Wein.

\spellinfobox{1}{1}{0 Schritt}{---}{---}{1 Kampfhandlung}{Nein}


\subsection{Flatulentio}

\spellheader{Illusion}{Natur}{1 Kampfhandlung}
Der Zauberer erzeugt die Illusion, dass ein Lebewesen, das sich in einer Entfernung von höchstens 50 Schritten befindet, einen außergewöhnlich lauten, weithin hörbaren Furz von sich gibt. Neben dem offensichtlichen Geräusch breitet sich auch eine markante Duftwolke von der Person weg aus.

Je höher die \textbf{Stärke} des Zaubers und das \textbf{Magieniveau} ist, desto auffälliger ist die Flatulenz.

\begin{pquote}{Martin Gerhard Reisenberg}
Auch das menschliche Hinterteil verliert dann und wann den Status der rückwärtigen Dienste und wird zur Offensivwaffe.
\end{pquote}

\spellinfobox{2}{1}{50 Schritt}{---}{Wolke}{1 Kampfhandlung}{Nein}


\subsection{Levitar}

\spellheader{Kontrolle}{Energie}{1 Kampfhandlung}
Mit einer einfachen Geste entfesselt der Zaubernde eine unsichtbare Kraft, die Gegenstände durch die Luft bewegt. Levitar kann genutzt werden, um Objekte aufzuheben, heranzuziehen oder mit einem kraftvollen Stoß fortzuschleudern. Die Stärke der Bewegung hängt von der Willenskraft des Anwenders ab.

Besonderheit: Erfahrene Magier können den Zauber verfeinern, um Objekte vorsichtig zu greifen oder sie im Kampf einzusetzen.

Spruchformeln: „Levitaris Volantis!“ – (Für präzise, schwebende Bewegungen) „Levitar Impetus!“ – (Für kräftige Stöße oder Würfe)

\spellinfobox{3}{1}{50 Schritt}{---}{---}{1 Kampfhandlung}{Nein}


\subsection{Gestaltwandler}

\spellheader{Illusion}{Licht}{1 Kampfhandlung}
Der Zaubernde verwandelt sich in eine Menschen ähnliche Gestalt. Es handelt sich dabei rein um die Illusion, Fähigkeiten werden nicht adaptiert.

\spellinfobox{5}{2}{0 Schritt}{---}{---}{1 Kampfhandlung}{Nein}


\subsection{Omniglossia (Der universelle Philologe)}

\spellheader{Verzauberung}{Arkana}{1 Kampfhandlung}
Wirkungsweise Dieser Zauber legt ein komplexes Netz aus magischer Energie über das Sprachzentrum des Wirkenden. Anstatt jede Sprache einzeln zu lernen, zapft der Magier das "Konzept der Sprache" an sich an.

Verständnis: Jede gehörte Sprache wird in Echtzeit im Bewusstsein in die Muttersprache übersetzt.

Artikulation: Der Zaubernde kann jede Sprache akzentfrei und fließend sprechen, als wäre er damit aufgewachsen.

Skriptorium: Unbekannte Schriftzeichen (sogar Hieroglyphen oder Runen) werden visuell so interpretiert, dass ihr Sinn klar vor dem geistigen Auge erscheint. Der Magier kann diese Zeichen ebenso flüssig verfassen.

\spellinfobox{5}{4}{0 Schritt}{---}{---}{1 Kampfhandlung}{Nein}


\subsection{Metamorphose der Urform}

\spellheader{Transmutation}{Arkana}{Ritual}
Wirkungsweise Dieser Zauber erlaubt es dem Wirkenden, die physische Form – inklusive Knochenstruktur, Gewebe und Organfunktionen – vollkommen neu zu ordnen. Die Transformation ist so fundamental, dass das Ziel nicht nur so aussieht wie die neue Gestalt, sondern deren biologische Eigenschaften (Stärke, Sinne, Fortbewegungsart) vollständig übernimmt.

Selbstverwandlung: Der Magier kann sich in jede natürliche Kreatur, ein anderes humanoides Wesen oder sogar in eine spezifische Person verwandeln.

Fremdverwandlung: Der Zauber kann auf ein Ziel in Reichweite gewirkt werden. Dies kann als Segen (Heilung von Deformationen, Anpassung an Umgebungen, etc.) oder als Fluch (Verwandlung eines Gegners in ein harmloses Tier) genutzt werden.

Teilaspekt-Wandlung: Es ist möglich, nur Teile des Körpers zu verändern (z. B. Kiemen wachsen lassen, die Haut verhärten oder Flügel formen).

\spellinfobox{5}{4}{20 Schritt}{---}{Sphäre}{Ritual}{Nein}


\end{multicols*}
